\documentclass{article}

\usepackage{amsmath, mathtools}

% Force the tables to appear where they should
\usepackage{float}

% Ensure that we can break URLs on a hyphen to avoid underfull hbox
% warnings in the bibliography.
\usepackage[hyphens]{url}

\usepackage{hyperref}
\hypersetup{
  colorlinks=true,
  linkcolor=blue,
  citecolor=blue
}

% The cone rank/dim table.
\usepackage{booktabs}
\renewcommand{\arraystretch}{1.5} % for the vertical padding
\setlength{\tabcolsep}{1em} % for the horizontal padding

%
% Abstract algebraic structures.
%
\ifx\havemjoalgebra\undefined
\def\havemjoalgebra{1}


\ifx\operatorname\undefined
  \usepackage{amsopn}
\fi

%
% Only the most commonly-used macros. Needed by everything else.
%
\ifx\havemjocommon\undefined
\def\havemjocommon{1}

\ifx\mathbb\undefined
  \usepackage{amsfonts}
\fi

\ifx\restriction\undefined
  \usepackage{amssymb}
\fi

% Place the argument in matching left/right parentheses.
\newcommand*{\of}[1]{ \left({#1}\right) }

% Group terms using parentheses.
\newcommand*{\qty}[1]{ \left({#1}\right) }

% Group terms using square brackets.
\newcommand*{\sqty}[1]{ \left[{#1}\right] }

% A pair of things.
\newcommand*{\pair}[2]{ \left({#1},{#2}\right) }

% A triple of things.
\newcommand*{\triple}[3]{ \left({#1},{#2},{#3}\right) }

% A four-tuple of things.
\newcommand*{\quadruple}[4]{ \left({#1},{#2},{#3},{#4}\right) }

% A five-tuple of things.
\newcommand*{\quintuple}[5]{ \left({#1},{#2},{#3},{#4},{#5}\right) }

% A six-tuple of things.
\newcommand*{\sextuple}[6]{ \left({#1},{#2},{#3},{#4},{#5},{#6}\right) }

% A seven-tuple of things.
\newcommand*{\septuple}[7]{ \left({#1},{#2},{#3},{#4},{#5},{#6},{#7}\right) }

% A free-form tuple of things. Useful for when the exact number is not
% known, such as when \ldots will be stuck in the middle of the list,
% and when you don't want to think in column-vector terms, e.g. with
% elements of an abstract Cartesian product space.
\newcommand*{\tuple}[1]{ \left({#1}\right) }

% The "least common multiple of" function. Takes a nonempty set of
% things that can be multiplied and ordered as its argument. Name
% chosen for synergy with \gcd, which *does* exist already.
\newcommand*{\lcm}[1]{ \operatorname{lcm}\of{{#1}} }
\ifdefined\newglossaryentry
  \newglossaryentry{lcm}{
    name={\ensuremath{\lcm{X}}},
    description={the least common multiple of the elements of $X$},
    sort=l
  }
\fi

% The factorial operator.
\newcommand*{\factorial}[1]{ {#1}! }

% Restrict the first argument (a function) to the second argument (a
% subset of that functions domain). Abused for polynomials to specify
% an associated function with a particular domain (also its codomain,
% in the case of univariate polynomials).
\newcommand*{\restrict}[2]{{#1}{\restriction}_{#2}}
\ifdefined\newglossaryentry
  \newglossaryentry{restriction}{
    name={\ensuremath{\restrict{f}{X}}},
    description={the restriction of $f$ to $X$},
    sort=r
  }
\fi

%
% Product spaces
%
% All of the product spaces (for example, R^n) that follow default to
% an exponent of ``n'', but that exponent can be changed by providing
% it as an optional argument. If the exponent given is ``1'', then it
% will be omitted entirely.
%

% The natural n-space, N x N x N x ... x N.
\newcommand*{\Nn}[1][n]{
  \mathbb{N}\if\detokenize{#1}\detokenize{1}{}\else^{#1}\fi
}

\ifdefined\newglossaryentry
  \newglossaryentry{N}{
    name={\ensuremath{\Nn[1]}},
    description={the set of natural numbers},
    sort=N
  }
\fi

% The integral n-space, Z x Z x Z x ... x Z.
\newcommand*{\Zn}[1][n]{
  \mathbb{Z}\if\detokenize{#1}\detokenize{1}{}\else^{#1}\fi
}

\ifdefined\newglossaryentry
  \newglossaryentry{Z}{
    name={\ensuremath{\Zn[1]}},
    description={the ring of integers},
    sort=Z
  }
\fi

% The rational n-space, Q x Q x Q x ... x Q.
\newcommand*{\Qn}[1][n]{
  \mathbb{Q}\if\detokenize{#1}\detokenize{1}{}\else^{#1}\fi
}

\ifdefined\newglossaryentry
  \newglossaryentry{Q}{
    name={\ensuremath{\Qn[1]}},
    description={the field of rational numbers},
    sort=Q
  }
\fi

% The real n-space, R x R x R x ... x R.
\newcommand*{\Rn}[1][n]{
  \mathbb{R}\if\detokenize{#1}\detokenize{1}{}\else^{#1}\fi
}

\ifdefined\newglossaryentry
  \newglossaryentry{R}{
    name={\ensuremath{\Rn[1]}},
    description={the field of real numbers},
    sort=R
  }
\fi


% The complex n-space, C x C x C x ... x C.
\newcommand*{\Cn}[1][n]{
  \mathbb{C}\if\detokenize{#1}\detokenize{1}{}\else^{#1}\fi
}

\ifdefined\newglossaryentry
  \newglossaryentry{C}{
    name={\ensuremath{\Cn[1]}},
    description={the field of complex numbers},
    sort=C
  }
\fi

% The n-dimensional product space of a generic field F.
\newcommand*{\Fn}[1][n]{
  \mathbb{F}\if\detokenize{#1}\detokenize{1}{}\else^{#1}\fi
}

\ifdefined\newglossaryentry
  \newglossaryentry{F}{
    name={\ensuremath{\Fn[1]}},
    description={a generic field},
    sort=F
  }
\fi


% An indexed arbitrary binary operation such as the union or
% intersection of an infinite number of sets. The first argument is
% the operator symbol to use, such as \cup for a union. The second
% argument is the lower index, for example k=1. The third argument is
% the upper index, such as \infty. Finally the fourth argument should
% contain the things (e.g. indexed sets) to be operated on.
\newcommand*{\binopmany}[4]{
  \mathchoice{ \underset{#2}{\overset{#3}{#1}}{#4} }
             { {#1}_{#2}^{#3}{#4} }
             { {#1}_{#2}^{#3}{#4} }
             { {#1}_{#2}^{#3}{#4} }
}


% The four standard (UNLESS YOU'RE FRENCH) types of intervals along
% the real line.
\newcommand*{\intervaloo}[2]{ \left({#1},{#2}\right) } % open-open
\newcommand*{\intervaloc}[2]{ \left({#1},{#2}\right] } % open-closed
\newcommand*{\intervalco}[2]{ \left[{#1},{#2}\right) } % closed-open
\newcommand*{\intervalcc}[2]{ \left[{#1},{#2}\right] } % closed-closed


\fi
 % for \of, and \binopmany


% The additive identity element of its argument, which should be
% an algebraic structure.
\newcommand*{\zero}[1]{ 0_{{#1}} }

\ifdefined\newglossaryentry
  \newglossaryentry{zero}{
    name={\ensuremath{\zero{R}}},
    description={the additive identity element of $R$},
    sort=z
  }
\fi

% The multiplicative identity element of its argument, which should be
% an algebraic structure.
\newcommand*{\unit}[1]{ 1_{{#1}} }

\ifdefined\newglossaryentry
  \newglossaryentry{unit}{
    name={\ensuremath{\unit{R}}},
    description={the multiplicative identity (unit) element of $R$},
    sort=u
  }
\fi

% The direct sum of two things.
\newcommand*{\directsum}[2]{ {#1}\oplus{#2} }

% The direct sum of three things.
\newcommand*{\directsumthree}[3]{ \directsum{#1}{\directsum{#2}{#3}} }

% The (indexed) direct sum of many things.
\newcommand*{\directsummany}[3]{ \binopmany{\bigoplus}{#1}{#2}{#3} }


% The (sub)algebra generated by its argument, a subset of some ambient
% algebra. By definition this is the smallest subalgebra (of the
% ambient one) containing that set.
\newcommand*{\alg}[1]{\operatorname{alg}\of{{#1}}}
\ifdefined\newglossaryentry
  \newglossaryentry{alg}{
    name={\ensuremath{\alg{X}}},
    description={the (sub)algebra generated by $X$},
    sort=a
  }
\fi


% The fraction field of its argument, an integral domain. The name
% "Frac" was chosen here instead of "Quot" because the latter
% corresponds to the term "quotient field," which can be mistaken in
% some cases for... a quotient field (something mod something).
\newcommand*{\Frac}[1]{\operatorname{Frac}\of{{#1}}}

% The ideal generated by its argument, a subset consisting of ring or
% algebra elements.
\newcommand*{\ideal}[1]{\operatorname{ideal}\of{{#1}}}
\ifdefined\newglossaryentry
  \newglossaryentry{ideal}{
    name={\ensuremath{\ideal{X}}},
    description={the ideal generated by $X$},
    sort=i
  }
\fi


% The polynomial ring whose underlying commutative ring of
% coefficients is the first argument and whose indeterminates (a
% comma-separated list) are the second argumnt.
\newcommand*{\polyring}[2]{{#1}\left[{#2}\right]}
\ifdefined\newglossaryentry
  \newglossaryentry{polyring}{
    name={\ensuremath{\polyring{R}{X}}},
    description={polynomials with coefficients in $R$ and variable $X$},
    sort=p
  }
\fi


% The stabilizer subgroup of its first argument that fixes the point
% given by its second argument.
\newcommand*{\Stab}[2]{ #1_{#2} }


% The affine algebraic variety consisting of the common solutions to
% every polynomial in its argument, which should be a subset of some
% polynomial ring.
\newcommand*{\variety}[1]{ \mathcal{V}\of{{#1}} }
\ifdefined\newglossaryentry
  \newglossaryentry{variety}{
    name={\ensuremath{\variety{I}}},
    description={variety corresponding to the ideal $I$},
    sort=v
  }
\fi

\fi

%
% Only the most commonly-used macros. Needed by everything else.
%
\ifx\havemjocommon\undefined
\def\havemjocommon{1}

\ifx\mathbb\undefined
  \usepackage{amsfonts}
\fi

\ifx\restriction\undefined
  \usepackage{amssymb}
\fi

% Place the argument in matching left/right parentheses.
\newcommand*{\of}[1]{ \left({#1}\right) }

% Group terms using parentheses.
\newcommand*{\qty}[1]{ \left({#1}\right) }

% Group terms using square brackets.
\newcommand*{\sqty}[1]{ \left[{#1}\right] }

% A pair of things.
\newcommand*{\pair}[2]{ \left({#1},{#2}\right) }

% A triple of things.
\newcommand*{\triple}[3]{ \left({#1},{#2},{#3}\right) }

% A four-tuple of things.
\newcommand*{\quadruple}[4]{ \left({#1},{#2},{#3},{#4}\right) }

% A five-tuple of things.
\newcommand*{\quintuple}[5]{ \left({#1},{#2},{#3},{#4},{#5}\right) }

% A six-tuple of things.
\newcommand*{\sextuple}[6]{ \left({#1},{#2},{#3},{#4},{#5},{#6}\right) }

% A seven-tuple of things.
\newcommand*{\septuple}[7]{ \left({#1},{#2},{#3},{#4},{#5},{#6},{#7}\right) }

% A free-form tuple of things. Useful for when the exact number is not
% known, such as when \ldots will be stuck in the middle of the list,
% and when you don't want to think in column-vector terms, e.g. with
% elements of an abstract Cartesian product space.
\newcommand*{\tuple}[1]{ \left({#1}\right) }

% The "least common multiple of" function. Takes a nonempty set of
% things that can be multiplied and ordered as its argument. Name
% chosen for synergy with \gcd, which *does* exist already.
\newcommand*{\lcm}[1]{ \operatorname{lcm}\of{{#1}} }
\ifdefined\newglossaryentry
  \newglossaryentry{lcm}{
    name={\ensuremath{\lcm{X}}},
    description={the least common multiple of the elements of $X$},
    sort=l
  }
\fi

% The factorial operator.
\newcommand*{\factorial}[1]{ {#1}! }

% Restrict the first argument (a function) to the second argument (a
% subset of that functions domain). Abused for polynomials to specify
% an associated function with a particular domain (also its codomain,
% in the case of univariate polynomials).
\newcommand*{\restrict}[2]{{#1}{\restriction}_{#2}}
\ifdefined\newglossaryentry
  \newglossaryentry{restriction}{
    name={\ensuremath{\restrict{f}{X}}},
    description={the restriction of $f$ to $X$},
    sort=r
  }
\fi

%
% Product spaces
%
% All of the product spaces (for example, R^n) that follow default to
% an exponent of ``n'', but that exponent can be changed by providing
% it as an optional argument. If the exponent given is ``1'', then it
% will be omitted entirely.
%

% The natural n-space, N x N x N x ... x N.
\newcommand*{\Nn}[1][n]{
  \mathbb{N}\if\detokenize{#1}\detokenize{1}{}\else^{#1}\fi
}

\ifdefined\newglossaryentry
  \newglossaryentry{N}{
    name={\ensuremath{\Nn[1]}},
    description={the set of natural numbers},
    sort=N
  }
\fi

% The integral n-space, Z x Z x Z x ... x Z.
\newcommand*{\Zn}[1][n]{
  \mathbb{Z}\if\detokenize{#1}\detokenize{1}{}\else^{#1}\fi
}

\ifdefined\newglossaryentry
  \newglossaryentry{Z}{
    name={\ensuremath{\Zn[1]}},
    description={the ring of integers},
    sort=Z
  }
\fi

% The rational n-space, Q x Q x Q x ... x Q.
\newcommand*{\Qn}[1][n]{
  \mathbb{Q}\if\detokenize{#1}\detokenize{1}{}\else^{#1}\fi
}

\ifdefined\newglossaryentry
  \newglossaryentry{Q}{
    name={\ensuremath{\Qn[1]}},
    description={the field of rational numbers},
    sort=Q
  }
\fi

% The real n-space, R x R x R x ... x R.
\newcommand*{\Rn}[1][n]{
  \mathbb{R}\if\detokenize{#1}\detokenize{1}{}\else^{#1}\fi
}

\ifdefined\newglossaryentry
  \newglossaryentry{R}{
    name={\ensuremath{\Rn[1]}},
    description={the field of real numbers},
    sort=R
  }
\fi


% The complex n-space, C x C x C x ... x C.
\newcommand*{\Cn}[1][n]{
  \mathbb{C}\if\detokenize{#1}\detokenize{1}{}\else^{#1}\fi
}

\ifdefined\newglossaryentry
  \newglossaryentry{C}{
    name={\ensuremath{\Cn[1]}},
    description={the field of complex numbers},
    sort=C
  }
\fi

% The n-dimensional product space of a generic field F.
\newcommand*{\Fn}[1][n]{
  \mathbb{F}\if\detokenize{#1}\detokenize{1}{}\else^{#1}\fi
}

\ifdefined\newglossaryentry
  \newglossaryentry{F}{
    name={\ensuremath{\Fn[1]}},
    description={a generic field},
    sort=F
  }
\fi


% An indexed arbitrary binary operation such as the union or
% intersection of an infinite number of sets. The first argument is
% the operator symbol to use, such as \cup for a union. The second
% argument is the lower index, for example k=1. The third argument is
% the upper index, such as \infty. Finally the fourth argument should
% contain the things (e.g. indexed sets) to be operated on.
\newcommand*{\binopmany}[4]{
  \mathchoice{ \underset{#2}{\overset{#3}{#1}}{#4} }
             { {#1}_{#2}^{#3}{#4} }
             { {#1}_{#2}^{#3}{#4} }
             { {#1}_{#2}^{#3}{#4} }
}


% The four standard (UNLESS YOU'RE FRENCH) types of intervals along
% the real line.
\newcommand*{\intervaloo}[2]{ \left({#1},{#2}\right) } % open-open
\newcommand*{\intervaloc}[2]{ \left({#1},{#2}\right] } % open-closed
\newcommand*{\intervalco}[2]{ \left[{#1},{#2}\right) } % closed-open
\newcommand*{\intervalcc}[2]{ \left[{#1},{#2}\right] } % closed-closed


\fi

%
% Cone stuff.
%
% The operator families Z(K), LL(K), etc. can technically be defined on
% sets other than cones, but nobody cares.
%
\ifx\havemjocone\undefined
\def\havemjocone{1}


\ifx\succcurlyeq\undefined
  \usepackage{amssymb} % \succcurlyeq, \preccurlyeq
\fi

%
% Only the most commonly-used macros. Needed by everything else.
%
\ifx\havemjocommon\undefined
\def\havemjocommon{1}

\ifx\mathbb\undefined
  \usepackage{amsfonts}
\fi

\ifx\restriction\undefined
  \usepackage{amssymb}
\fi

% Place the argument in matching left/right parentheses.
\newcommand*{\of}[1]{ \left({#1}\right) }

% Group terms using parentheses.
\newcommand*{\qty}[1]{ \left({#1}\right) }

% Group terms using square brackets.
\newcommand*{\sqty}[1]{ \left[{#1}\right] }

% A pair of things.
\newcommand*{\pair}[2]{ \left({#1},{#2}\right) }

% A triple of things.
\newcommand*{\triple}[3]{ \left({#1},{#2},{#3}\right) }

% A four-tuple of things.
\newcommand*{\quadruple}[4]{ \left({#1},{#2},{#3},{#4}\right) }

% A five-tuple of things.
\newcommand*{\quintuple}[5]{ \left({#1},{#2},{#3},{#4},{#5}\right) }

% A six-tuple of things.
\newcommand*{\sextuple}[6]{ \left({#1},{#2},{#3},{#4},{#5},{#6}\right) }

% A seven-tuple of things.
\newcommand*{\septuple}[7]{ \left({#1},{#2},{#3},{#4},{#5},{#6},{#7}\right) }

% A free-form tuple of things. Useful for when the exact number is not
% known, such as when \ldots will be stuck in the middle of the list,
% and when you don't want to think in column-vector terms, e.g. with
% elements of an abstract Cartesian product space.
\newcommand*{\tuple}[1]{ \left({#1}\right) }

% The "least common multiple of" function. Takes a nonempty set of
% things that can be multiplied and ordered as its argument. Name
% chosen for synergy with \gcd, which *does* exist already.
\newcommand*{\lcm}[1]{ \operatorname{lcm}\of{{#1}} }
\ifdefined\newglossaryentry
  \newglossaryentry{lcm}{
    name={\ensuremath{\lcm{X}}},
    description={the least common multiple of the elements of $X$},
    sort=l
  }
\fi

% The factorial operator.
\newcommand*{\factorial}[1]{ {#1}! }

% Restrict the first argument (a function) to the second argument (a
% subset of that functions domain). Abused for polynomials to specify
% an associated function with a particular domain (also its codomain,
% in the case of univariate polynomials).
\newcommand*{\restrict}[2]{{#1}{\restriction}_{#2}}
\ifdefined\newglossaryentry
  \newglossaryentry{restriction}{
    name={\ensuremath{\restrict{f}{X}}},
    description={the restriction of $f$ to $X$},
    sort=r
  }
\fi

%
% Product spaces
%
% All of the product spaces (for example, R^n) that follow default to
% an exponent of ``n'', but that exponent can be changed by providing
% it as an optional argument. If the exponent given is ``1'', then it
% will be omitted entirely.
%

% The natural n-space, N x N x N x ... x N.
\newcommand*{\Nn}[1][n]{
  \mathbb{N}\if\detokenize{#1}\detokenize{1}{}\else^{#1}\fi
}

\ifdefined\newglossaryentry
  \newglossaryentry{N}{
    name={\ensuremath{\Nn[1]}},
    description={the set of natural numbers},
    sort=N
  }
\fi

% The integral n-space, Z x Z x Z x ... x Z.
\newcommand*{\Zn}[1][n]{
  \mathbb{Z}\if\detokenize{#1}\detokenize{1}{}\else^{#1}\fi
}

\ifdefined\newglossaryentry
  \newglossaryentry{Z}{
    name={\ensuremath{\Zn[1]}},
    description={the ring of integers},
    sort=Z
  }
\fi

% The rational n-space, Q x Q x Q x ... x Q.
\newcommand*{\Qn}[1][n]{
  \mathbb{Q}\if\detokenize{#1}\detokenize{1}{}\else^{#1}\fi
}

\ifdefined\newglossaryentry
  \newglossaryentry{Q}{
    name={\ensuremath{\Qn[1]}},
    description={the field of rational numbers},
    sort=Q
  }
\fi

% The real n-space, R x R x R x ... x R.
\newcommand*{\Rn}[1][n]{
  \mathbb{R}\if\detokenize{#1}\detokenize{1}{}\else^{#1}\fi
}

\ifdefined\newglossaryentry
  \newglossaryentry{R}{
    name={\ensuremath{\Rn[1]}},
    description={the field of real numbers},
    sort=R
  }
\fi


% The complex n-space, C x C x C x ... x C.
\newcommand*{\Cn}[1][n]{
  \mathbb{C}\if\detokenize{#1}\detokenize{1}{}\else^{#1}\fi
}

\ifdefined\newglossaryentry
  \newglossaryentry{C}{
    name={\ensuremath{\Cn[1]}},
    description={the field of complex numbers},
    sort=C
  }
\fi

% The n-dimensional product space of a generic field F.
\newcommand*{\Fn}[1][n]{
  \mathbb{F}\if\detokenize{#1}\detokenize{1}{}\else^{#1}\fi
}

\ifdefined\newglossaryentry
  \newglossaryentry{F}{
    name={\ensuremath{\Fn[1]}},
    description={a generic field},
    sort=F
  }
\fi


% An indexed arbitrary binary operation such as the union or
% intersection of an infinite number of sets. The first argument is
% the operator symbol to use, such as \cup for a union. The second
% argument is the lower index, for example k=1. The third argument is
% the upper index, such as \infty. Finally the fourth argument should
% contain the things (e.g. indexed sets) to be operated on.
\newcommand*{\binopmany}[4]{
  \mathchoice{ \underset{#2}{\overset{#3}{#1}}{#4} }
             { {#1}_{#2}^{#3}{#4} }
             { {#1}_{#2}^{#3}{#4} }
             { {#1}_{#2}^{#3}{#4} }
}


% The four standard (UNLESS YOU'RE FRENCH) types of intervals along
% the real line.
\newcommand*{\intervaloo}[2]{ \left({#1},{#2}\right) } % open-open
\newcommand*{\intervaloc}[2]{ \left({#1},{#2}\right] } % open-closed
\newcommand*{\intervalco}[2]{ \left[{#1},{#2}\right) } % closed-open
\newcommand*{\intervalcc}[2]{ \left[{#1},{#2}\right] } % closed-closed


\fi
 % for \of, \Rn, etc.
%
% Standard operations from linear algebra.
%
\ifx\havemjolinearalgebra\undefined
\def\havemjolinearalgebra{1}


\ifx\lvert\undefined
  \usepackage{amsmath} % \lvert, \rVert, etc. and \operatorname.
\fi

\ifx\ocircle\undefined
  \usepackage{wasysym}
\fi

\ifx\clipbox\undefined
  % Part of the adjustbox package; needed to clip the \perp sign.
  \usepackage{trimclip}
\fi

%
% Only the most commonly-used macros. Needed by everything else.
%
\ifx\havemjocommon\undefined
\def\havemjocommon{1}

\ifx\mathbb\undefined
  \usepackage{amsfonts}
\fi

\ifx\restriction\undefined
  \usepackage{amssymb}
\fi

% Place the argument in matching left/right parentheses.
\newcommand*{\of}[1]{ \left({#1}\right) }

% Group terms using parentheses.
\newcommand*{\qty}[1]{ \left({#1}\right) }

% Group terms using square brackets.
\newcommand*{\sqty}[1]{ \left[{#1}\right] }

% A pair of things.
\newcommand*{\pair}[2]{ \left({#1},{#2}\right) }

% A triple of things.
\newcommand*{\triple}[3]{ \left({#1},{#2},{#3}\right) }

% A four-tuple of things.
\newcommand*{\quadruple}[4]{ \left({#1},{#2},{#3},{#4}\right) }

% A five-tuple of things.
\newcommand*{\quintuple}[5]{ \left({#1},{#2},{#3},{#4},{#5}\right) }

% A six-tuple of things.
\newcommand*{\sextuple}[6]{ \left({#1},{#2},{#3},{#4},{#5},{#6}\right) }

% A seven-tuple of things.
\newcommand*{\septuple}[7]{ \left({#1},{#2},{#3},{#4},{#5},{#6},{#7}\right) }

% A free-form tuple of things. Useful for when the exact number is not
% known, such as when \ldots will be stuck in the middle of the list,
% and when you don't want to think in column-vector terms, e.g. with
% elements of an abstract Cartesian product space.
\newcommand*{\tuple}[1]{ \left({#1}\right) }

% The "least common multiple of" function. Takes a nonempty set of
% things that can be multiplied and ordered as its argument. Name
% chosen for synergy with \gcd, which *does* exist already.
\newcommand*{\lcm}[1]{ \operatorname{lcm}\of{{#1}} }
\ifdefined\newglossaryentry
  \newglossaryentry{lcm}{
    name={\ensuremath{\lcm{X}}},
    description={the least common multiple of the elements of $X$},
    sort=l
  }
\fi

% The factorial operator.
\newcommand*{\factorial}[1]{ {#1}! }

% Restrict the first argument (a function) to the second argument (a
% subset of that functions domain). Abused for polynomials to specify
% an associated function with a particular domain (also its codomain,
% in the case of univariate polynomials).
\newcommand*{\restrict}[2]{{#1}{\restriction}_{#2}}
\ifdefined\newglossaryentry
  \newglossaryentry{restriction}{
    name={\ensuremath{\restrict{f}{X}}},
    description={the restriction of $f$ to $X$},
    sort=r
  }
\fi

%
% Product spaces
%
% All of the product spaces (for example, R^n) that follow default to
% an exponent of ``n'', but that exponent can be changed by providing
% it as an optional argument. If the exponent given is ``1'', then it
% will be omitted entirely.
%

% The natural n-space, N x N x N x ... x N.
\newcommand*{\Nn}[1][n]{
  \mathbb{N}\if\detokenize{#1}\detokenize{1}{}\else^{#1}\fi
}

\ifdefined\newglossaryentry
  \newglossaryentry{N}{
    name={\ensuremath{\Nn[1]}},
    description={the set of natural numbers},
    sort=N
  }
\fi

% The integral n-space, Z x Z x Z x ... x Z.
\newcommand*{\Zn}[1][n]{
  \mathbb{Z}\if\detokenize{#1}\detokenize{1}{}\else^{#1}\fi
}

\ifdefined\newglossaryentry
  \newglossaryentry{Z}{
    name={\ensuremath{\Zn[1]}},
    description={the ring of integers},
    sort=Z
  }
\fi

% The rational n-space, Q x Q x Q x ... x Q.
\newcommand*{\Qn}[1][n]{
  \mathbb{Q}\if\detokenize{#1}\detokenize{1}{}\else^{#1}\fi
}

\ifdefined\newglossaryentry
  \newglossaryentry{Q}{
    name={\ensuremath{\Qn[1]}},
    description={the field of rational numbers},
    sort=Q
  }
\fi

% The real n-space, R x R x R x ... x R.
\newcommand*{\Rn}[1][n]{
  \mathbb{R}\if\detokenize{#1}\detokenize{1}{}\else^{#1}\fi
}

\ifdefined\newglossaryentry
  \newglossaryentry{R}{
    name={\ensuremath{\Rn[1]}},
    description={the field of real numbers},
    sort=R
  }
\fi


% The complex n-space, C x C x C x ... x C.
\newcommand*{\Cn}[1][n]{
  \mathbb{C}\if\detokenize{#1}\detokenize{1}{}\else^{#1}\fi
}

\ifdefined\newglossaryentry
  \newglossaryentry{C}{
    name={\ensuremath{\Cn[1]}},
    description={the field of complex numbers},
    sort=C
  }
\fi

% The n-dimensional product space of a generic field F.
\newcommand*{\Fn}[1][n]{
  \mathbb{F}\if\detokenize{#1}\detokenize{1}{}\else^{#1}\fi
}

\ifdefined\newglossaryentry
  \newglossaryentry{F}{
    name={\ensuremath{\Fn[1]}},
    description={a generic field},
    sort=F
  }
\fi


% An indexed arbitrary binary operation such as the union or
% intersection of an infinite number of sets. The first argument is
% the operator symbol to use, such as \cup for a union. The second
% argument is the lower index, for example k=1. The third argument is
% the upper index, such as \infty. Finally the fourth argument should
% contain the things (e.g. indexed sets) to be operated on.
\newcommand*{\binopmany}[4]{
  \mathchoice{ \underset{#2}{\overset{#3}{#1}}{#4} }
             { {#1}_{#2}^{#3}{#4} }
             { {#1}_{#2}^{#3}{#4} }
             { {#1}_{#2}^{#3}{#4} }
}


% The four standard (UNLESS YOU'RE FRENCH) types of intervals along
% the real line.
\newcommand*{\intervaloo}[2]{ \left({#1},{#2}\right) } % open-open
\newcommand*{\intervaloc}[2]{ \left({#1},{#2}\right] } % open-closed
\newcommand*{\intervalco}[2]{ \left[{#1},{#2}\right) } % closed-open
\newcommand*{\intervalcc}[2]{ \left[{#1},{#2}\right] } % closed-closed


\fi
 % for \of, at least

% Absolute value (modulus) of a scalar.
\newcommand*{\abs}[1]{\left\lvert{#1}\right\rvert}

% Norm of a vector.
\newcommand*{\norm}[1]{\left\lVert{#1}\right\rVert}

% The inner product between its two arguments.
\newcommand*{\ip}[2]{\left\langle{#1},{#2}\right\rangle}

% The tensor product of its two arguments.
\newcommand*{\tp}[2]{ {#1}\otimes{#2} }

% The Kronecker product of its two arguments. The usual notation for
% this is the same as the tensor product notation used for \tp, but
% that leads to confusion because the two definitions may not agree.
\newcommand*{\kp}[2]{ {#1}\odot{#2} }

% The adjoint of a linear operator.
\newcommand*{\adjoint}[1]{ #1^{*} }

% The ``transpose'' of a linear operator; namely, the adjoint, but
% specialized to real matrices.
\newcommand*{\transpose}[1]{ #1^{T} }

% The Moore-Penrose (or any other, I guess) pseudo-inverse of its
% sole argument.
\newcommand*{\pseudoinverse}[1]{ #1^{+} }

% The trace of an operator.
\newcommand*{\trace}[1]{ \operatorname{trace}\of{{#1}} }

% The diagonal matrix whose only nonzero entries are on the diagonal
% and are given by our argument. The argument should therefore be a
% vector or tuple of entries, by convention going from the top-left to
% the bottom-right of the matrix.
\newcommand*{\diag}[1]{\operatorname{diag}\of{{#1}}}

% The "rank" of its argument, which is context-dependent. It can mean
% any or all of,
%
%   * the rank of a matrix,
%   * the rank of a power-associative algebra (particularly an EJA),
%   * the rank of an element in a Euclidean Jordan algebra.
%
\newcommand*{\rank}[1]{ \operatorname{rank}\of{{#1}} }


% The ``span of'' operator. The name \span is already taken.
\newcommand*{\spanof}[1]{ \operatorname{span}\of{{#1}} }

% The ``co-dimension of'' operator.
\newcommand*{\codim}{ \operatorname{codim} }

% The orthogonal projection of its second argument onto the first.
\newcommand*{\proj}[2] { \operatorname{proj}\of{#1, #2} }

% The set of all eigenvalues of its argument, which should be either a
% matrix or a linear operator. The sigma notation was chosen instead
% of lambda so that lambda can be reserved to denote the ordered tuple
% (largest to smallest) of eigenvalues.
\newcommand*{\spectrum}[1]{\sigma\of{{#1}}}
\ifdefined\newglossaryentry
  \newglossaryentry{spectrum}{
    name={\ensuremath{\spectrum{L}}},
    description={the set of all eigenvalues of $L$},
    sort=s
  }
\fi

% The reduced row-echelon form of its argument, a matrix.
\newcommand*{\rref}[1]{\operatorname{rref}\of{#1}}
\ifdefined\newglossaryentry
  \newglossaryentry{rref}{
    name={\ensuremath{\rref{A}}},
    description={the reduced row-echelon form of $A$},
    sort=r
  }
\fi

% The ``Automorphism group of'' operator.
\newcommand*{\Aut}[1]{ \operatorname{Aut}\of{{#1}} }

% The ``Lie algebra of'' operator.
\newcommand*{\Lie}[1]{ \operatorname{Lie}\of{{#1}} }

% The ``write a matrix as a big vector'' operator.
\newcommand*{\vectorize}[1]{ \operatorname{vec}\of{{#1}} }

% The ``write a big vector as a matrix'' operator.
\newcommand*{\matricize}[1]{ \operatorname{mat}\of{{#1}} }

% An inline column vector, with parentheses and a transpose operator.
\newcommand*{\colvec}[1]{ \transpose{\left({#1}\right)} }

% Bounded linear operators on some space. The required argument is the
% domain of those operators, and the optional argument is the
% codomain. If the optional argument is omitted, the required argument
% is used for both.
\newcommand*{\boundedops}[2][]{
  \mathcal{B}\of{ {#2}
    \if\relax\detokenize{#1}\relax
      {}%
    \else
      {,{#1}}%
    \fi
  }
}


%
% Orthogonal direct sum.
%
% First declare my ``perp in a circle'' operator, which is meant to be
% like an \obot or an \operp except has the correct weight circle. It's
% achieved by overlaying an \ocircle with a \perp, but only after we
% clip off the top half of the \perp sign and shift it up.
\DeclareMathOperator{\oplusperp}{\mathbin{
  \ooalign{
    $\ocircle$\cr
    \raisebox{0.625\height}{$\clipbox{0pt 0pt 0pt 0.5\height}{$\perp$}$}\cr
  }
}}

% Now declare an orthogonal direct sum in terms of \oplusperp.
\newcommand*{\directsumperp}[2]{ {#1}\oplusperp{#2} }


% The space of real symmetric n-by-n matrices. Does not reduce to
% merely "S" when n=1 since S^{n} does not mean an n-fold cartesian
% product of S^{1}.
\newcommand*{\Sn}[1][n]{ \mathcal{S}^{#1} }
\ifdefined\newglossaryentry
  \newglossaryentry{Sn}{
    name={\ensuremath{\Sn}},
    description={the set of $n$-by-$n$ real symmetric matrices},
    sort=Sn
  }
\fi

% The space of complex Hermitian n-by-n matrices. Does not reduce to
% merely "H" when n=1 since H^{n} does not mean an n-fold cartesian
% product of H^{1}. The field may also be given rather than assumed
% to be complex; for example \Hn[3]\of{\mathbb{O}} might denote the
% 3-by-3 Hermitian matrices with octonion entries.
\newcommand*{\Hn}[1][n]{ \mathcal{H}^{#1} }
\ifdefined\newglossaryentry
  \newglossaryentry{Hn}{
    name={\ensuremath{\Hn}},
    description={the set of $n$-by-$n$ complex Hermitian matrices},
    sort=Hn
  }
\fi


% The general linear group of square matrices whose size is the first
% argument and whose entries come from the second argument.
\newcommand*{\GL}[2]{\operatorname{GL}_{#1}\of{#2}}
\ifdefined\newglossaryentry
  \newglossaryentry{GL}{
    name={\ensuremath{\GL{n}{\mathbb{F}}}},
    description={the general linear group of order $n$ over $\mathbb{F}$},
    sort=GL
  }
\fi


% The space of all (linear) derivations on its argument, an algebra.
\newcommand*{\Der}[1]{\operatorname{Der}\of{#1}}
\ifdefined\newglossaryentry
  \newglossaryentry{Der}{
    name={\ensuremath{\Der{V}}},
    description={the space of all (linear) derivations on $V$},
    sort=Der
  }
\fi


% The set of all isometries from its first argument to its second
% (optional) argument, both assumed to be at least normed spaces, and
% more likely Hilbert spaces. The norms are implicit, i.e. not
% included in the arguments. If the optional argument is omitted, you
% get the isometries from the first argument to itself AKA its
% isometry group.
\newcommand*{\Isom}[2][]{
  \operatorname{Isom}\of{ {#2}
    \if\relax\detokenize{#1}\relax
      {}%
    \else
      {,{#1}}%
    \fi
  }
}
\ifdefined\newglossaryentry
  \newglossaryentry{Isom}{
    name={\ensuremath{\Isom{V}}},
    description={the group of isometries on $V$},
    sort=Isom
  }
  \newglossaryentry{Isom2}{
    name={\ensuremath{\Isom[W]{V}}},
    description={the set of isometries from $V$ to $W$},
    sort=Isom
  }
\fi

\fi
 % \Sn and \Hn

% The dual of a subset of an inner-product space; always a closed
% convex cone.
\newcommand*{\dual}[1]{ #1^{*} }

%
% Common cones.
%

% The nonnegative and strictly positive orthants in the given number
% of dimensions.
\newcommand*{\Rnplus}[1][n]{ \Rn[#1]_{+} }
\newcommand*{\Rnplusplus}[1][n]{ \Rn[#1]_{++} }

% The Lorentz ``ice-cream'' cone in the given number of dimensions.
\newcommand*{\Lnplus}[1][n]{ \mathcal{L}^{{#1}}_{+} }

% The PSD cone in a space of symmetric matrices.
\newcommand*{\Snplus}[1][n]{ \Sn[#1]_{+} }

% The PSD cone in a space of Hermitian matrices.
\newcommand*{\Hnplus}[1][n]{ \Hn[#1]_{+} }


%
% Some collections of linear operators.
%

% The set of all positive operators on its argument. This uses the
% same magic as \boundedops to accept either one or two arguments. If
% one argument is given, the domain and codomain are equal and the
% positive operators fix a subset of that space. When two arguments
% are given, the positive operators send the first argument to a
% subset of the second.
\newcommand*{\posops}[2][]{
  \pi\of{ {#2}
    \if\relax\detokenize{#1}\relax
      {}%
    \else
      {,{#1}}%
    \fi
  }
}

% The set of all S-operators on its argument.
\newcommand*{\Sof}[1]{ \mathbf{S} \of{ {#1} } }

% The cone of all Z-operators on its argument.
\newcommand*{\Zof}[1]{ \mathbf{Z} \of{ {#1} } }

% The space of Lyapunov-like operators on its argument.
\newcommand*{\LL}[1]{ \mathbf{LL}\of{ {#1} } }

% The Lyapunov rank of the given cone.
\newcommand*{\lyapunovrank}[1]{ \beta\of{ {#1} } }

% Cone inequality operators.
\newcommand*{\gecone}{\succcurlyeq}
\newcommand*{\gtcone}{\succ}
\newcommand*{\lecone}{\preccurlyeq}
\newcommand*{\ltcone}{\prec}


\fi

%
% Font handling that you probably want to use everywhere.
%
\ifx\havemjofont\undefined
\def\havemjofont{1}

% Support utf8 in LaTeX code. This lets you enter names like
% Carathéodory and Güler directly, without the funny escape sequences.
\usepackage[utf8]{inputenc}

% Grab the AMS fonts, too, for things like \mathbb.
\ifx\mathbb\undefined
  \usepackage{amsfonts}
\fi

% Add the \mathpzc (PostScript Zapf Chancery) font face, scaled up by
% a factor of 1.15. This is described in the documentation for the
% urwchancal package which further adapts this technique. So long
% as we don't ask too much of it, however, the few lines below suffice.
\DeclareFontFamily{OT1}{pzc}{}
\DeclareFontShape{OT1}{pzc}{m}{it}{<-> s * [1.15] pzcmi7t}{}
\DeclareMathAlphabet{\mathpzc}{OT1}{pzc}{m}{it}

\fi

\ifx\havemjohurwitz\undefined
\def\havemjohurwitz{1}


\newcommand*{\quaternions}{\mathbb{H}}

\ifdefined\newglossaryentry
  \newglossaryentry{quaternions}{
    name={\ensuremath{\quaternions}},
    description={the algebra of quaternions},
    sort=H
  }
\fi


\newcommand*{\octonions}{\mathbb{O}}

\ifdefined\newglossaryentry
  \newglossaryentry{octonions}{
    name={\ensuremath{\octonions}},
    description={the algebra of octonions},
    sort=O
  }
\fi


\fi

%
% A proof-by-cases environment. This gives you a nicely-indented list
% of cases for use in proofs.
%
% Example:
%
% Case 1 (x >= 0): herp.
%
% Case 2 (x < 0): derp.
%
\ifx\havemjoproofbycases\undefined
\def\havemjoproofbycases{1}


% Used below to define pcases. The ``loadonly'' parameter prevents
% a very bad interaction with the beamer document class.
\usepackage[loadonly]{enumitem}

% Needed to perform division in the definition of \singleblskip.
\usepackage{calc}

% A \baselineskip without the \baselinestretch scaling factor. Even
% though \baselinestretch defaults to 1.0, it doesn't really have that
% value. Thus to avoid division by zero, we need to do the ``is this
% thing empty?'' hack.
%
% If we use \baselineskip instead of \singleblskip in our list, things
% get real ugly when the text is e.g. double-spaced.
%
\newlength{\singleblskip}
\setlength{\singleblskip}{
  \baselineskip/\real{
      \if\relax\detokenize{\baselinestretch}\relax
        \baselinestretch%
      \else
        1%
      \fi}
}

% Using the enumitem package, we define a new type of list, called
% ``pcases'' (proof by cases). Each case has a label with an arabic
% numeral (the case number), but also a \thiscase identifier. The
% macro \thiscase is defined below by the \case command, and gives the
% name or conditions or whatever that distinguish one case from
% another.
\newlist{pcases}{enumerate}{1}
\setlist[pcases]{
  label=\textit{Case~\arabic*}:~\protect\thiscase\textit{.},
  ref=\arabic*,
  align=left,
  listparindent=\parindent,
  parsep=\parskip,
  itemsep=\singleblskip}

% The optional argument here gets stuffed into the \thiscase macro, to
% be called by pcases when it creates this list item. The \hfill is a
% hack intended to force the proof to start on a new line, rather than
% right after the colon. A \newline where the \hfill is does not work,
% so we consume the rest of the line instead.
\newcommand{\case}[1][]{
  \def\thiscase{#1}%
  \item \hfill\par\vspace{\singleblskip}
}


\fi

%
% Things that fit absolutely nowhere else.
%
\ifx\havemjoset\undefined
\def\havemjoset{1}

%
% Only the most commonly-used macros. Needed by everything else.
%
\ifx\havemjocommon\undefined
\def\havemjocommon{1}

\ifx\mathbb\undefined
  \usepackage{amsfonts}
\fi

\ifx\restriction\undefined
  \usepackage{amssymb}
\fi

% Place the argument in matching left/right parentheses.
\newcommand*{\of}[1]{ \left({#1}\right) }

% Group terms using parentheses.
\newcommand*{\qty}[1]{ \left({#1}\right) }

% Group terms using square brackets.
\newcommand*{\sqty}[1]{ \left[{#1}\right] }

% A pair of things.
\newcommand*{\pair}[2]{ \left({#1},{#2}\right) }

% A triple of things.
\newcommand*{\triple}[3]{ \left({#1},{#2},{#3}\right) }

% A four-tuple of things.
\newcommand*{\quadruple}[4]{ \left({#1},{#2},{#3},{#4}\right) }

% A five-tuple of things.
\newcommand*{\quintuple}[5]{ \left({#1},{#2},{#3},{#4},{#5}\right) }

% A six-tuple of things.
\newcommand*{\sextuple}[6]{ \left({#1},{#2},{#3},{#4},{#5},{#6}\right) }

% A seven-tuple of things.
\newcommand*{\septuple}[7]{ \left({#1},{#2},{#3},{#4},{#5},{#6},{#7}\right) }

% A free-form tuple of things. Useful for when the exact number is not
% known, such as when \ldots will be stuck in the middle of the list,
% and when you don't want to think in column-vector terms, e.g. with
% elements of an abstract Cartesian product space.
\newcommand*{\tuple}[1]{ \left({#1}\right) }

% The "least common multiple of" function. Takes a nonempty set of
% things that can be multiplied and ordered as its argument. Name
% chosen for synergy with \gcd, which *does* exist already.
\newcommand*{\lcm}[1]{ \operatorname{lcm}\of{{#1}} }
\ifdefined\newglossaryentry
  \newglossaryentry{lcm}{
    name={\ensuremath{\lcm{X}}},
    description={the least common multiple of the elements of $X$},
    sort=l
  }
\fi

% The factorial operator.
\newcommand*{\factorial}[1]{ {#1}! }

% Restrict the first argument (a function) to the second argument (a
% subset of that functions domain). Abused for polynomials to specify
% an associated function with a particular domain (also its codomain,
% in the case of univariate polynomials).
\newcommand*{\restrict}[2]{{#1}{\restriction}_{#2}}
\ifdefined\newglossaryentry
  \newglossaryentry{restriction}{
    name={\ensuremath{\restrict{f}{X}}},
    description={the restriction of $f$ to $X$},
    sort=r
  }
\fi

%
% Product spaces
%
% All of the product spaces (for example, R^n) that follow default to
% an exponent of ``n'', but that exponent can be changed by providing
% it as an optional argument. If the exponent given is ``1'', then it
% will be omitted entirely.
%

% The natural n-space, N x N x N x ... x N.
\newcommand*{\Nn}[1][n]{
  \mathbb{N}\if\detokenize{#1}\detokenize{1}{}\else^{#1}\fi
}

\ifdefined\newglossaryentry
  \newglossaryentry{N}{
    name={\ensuremath{\Nn[1]}},
    description={the set of natural numbers},
    sort=N
  }
\fi

% The integral n-space, Z x Z x Z x ... x Z.
\newcommand*{\Zn}[1][n]{
  \mathbb{Z}\if\detokenize{#1}\detokenize{1}{}\else^{#1}\fi
}

\ifdefined\newglossaryentry
  \newglossaryentry{Z}{
    name={\ensuremath{\Zn[1]}},
    description={the ring of integers},
    sort=Z
  }
\fi

% The rational n-space, Q x Q x Q x ... x Q.
\newcommand*{\Qn}[1][n]{
  \mathbb{Q}\if\detokenize{#1}\detokenize{1}{}\else^{#1}\fi
}

\ifdefined\newglossaryentry
  \newglossaryentry{Q}{
    name={\ensuremath{\Qn[1]}},
    description={the field of rational numbers},
    sort=Q
  }
\fi

% The real n-space, R x R x R x ... x R.
\newcommand*{\Rn}[1][n]{
  \mathbb{R}\if\detokenize{#1}\detokenize{1}{}\else^{#1}\fi
}

\ifdefined\newglossaryentry
  \newglossaryentry{R}{
    name={\ensuremath{\Rn[1]}},
    description={the field of real numbers},
    sort=R
  }
\fi


% The complex n-space, C x C x C x ... x C.
\newcommand*{\Cn}[1][n]{
  \mathbb{C}\if\detokenize{#1}\detokenize{1}{}\else^{#1}\fi
}

\ifdefined\newglossaryentry
  \newglossaryentry{C}{
    name={\ensuremath{\Cn[1]}},
    description={the field of complex numbers},
    sort=C
  }
\fi

% The n-dimensional product space of a generic field F.
\newcommand*{\Fn}[1][n]{
  \mathbb{F}\if\detokenize{#1}\detokenize{1}{}\else^{#1}\fi
}

\ifdefined\newglossaryentry
  \newglossaryentry{F}{
    name={\ensuremath{\Fn[1]}},
    description={a generic field},
    sort=F
  }
\fi


% An indexed arbitrary binary operation such as the union or
% intersection of an infinite number of sets. The first argument is
% the operator symbol to use, such as \cup for a union. The second
% argument is the lower index, for example k=1. The third argument is
% the upper index, such as \infty. Finally the fourth argument should
% contain the things (e.g. indexed sets) to be operated on.
\newcommand*{\binopmany}[4]{
  \mathchoice{ \underset{#2}{\overset{#3}{#1}}{#4} }
             { {#1}_{#2}^{#3}{#4} }
             { {#1}_{#2}^{#3}{#4} }
             { {#1}_{#2}^{#3}{#4} }
}


% The four standard (UNLESS YOU'RE FRENCH) types of intervals along
% the real line.
\newcommand*{\intervaloo}[2]{ \left({#1},{#2}\right) } % open-open
\newcommand*{\intervaloc}[2]{ \left({#1},{#2}\right] } % open-closed
\newcommand*{\intervalco}[2]{ \left[{#1},{#2}\right) } % closed-open
\newcommand*{\intervalcc}[2]{ \left[{#1},{#2}\right] } % closed-closed


\fi
 % binopmany
%
% Font handling that you probably want to use everywhere.
%
\ifx\havemjofont\undefined
\def\havemjofont{1}

% Support utf8 in LaTeX code. This lets you enter names like
% Carathéodory and Güler directly, without the funny escape sequences.
\usepackage[utf8]{inputenc}

% Grab the AMS fonts, too, for things like \mathbb.
\ifx\mathbb\undefined
  \usepackage{amsfonts}
\fi

% Add the \mathpzc (PostScript Zapf Chancery) font face, scaled up by
% a factor of 1.15. This is described in the documentation for the
% urwchancal package which further adapts this technique. So long
% as we don't ask too much of it, however, the few lines below suffice.
\DeclareFontFamily{OT1}{pzc}{}
\DeclareFontShape{OT1}{pzc}{m}{it}{<-> s * [1.15] pzcmi7t}{}
\DeclareMathAlphabet{\mathpzc}{OT1}{pzc}{m}{it}

\fi
 % amsfonts and \mathpzc

\ifx\operatorname\undefined
  \usepackage{amsopn}
\fi

\ifx\bigtimes\undefined
  \usepackage{mathtools}
\fi

% Create a set from the given elements
\newcommand*{\set}[1]{\left\lbrace{#1}\right\rbrace}

% A set comprehension, where the ``such that...'' bar is added
% automatically. The bar was chosen over a colon to avoid ambiguity
% with the L : V -> V notation. We can't leverage \set here because \middle
% needs \left and \right present.
\newcommand*{\setc}[2]{\left\lbrace{#1}\ \middle|\ {#2}\right\rbrace}


% The cardinality of a set. The |X| notation conflicts with the
% absolute value, and the meaning of card(X) is clear at once, so we
% prefer the latter.
\newcommand*{\card}[1]{ \operatorname{card}\of{{#1}} }


% The powerset of (that is, the set of all subsets of) its argument.
\newcommand*{\powerset}[1]{\mathpzc{P}\of{{#1}}}
\ifdefined\newglossaryentry
  \newglossaryentry{powerset}{
    name={\ensuremath{\powerset{X}}},
    description={the ``powerset,'' or set of all subsets of $X$},
    sort=p
  }
\fi


%
% Basic set operations
%

% The union of its two arguments.
\newcommand*{\union}[2]{ {#1}\cup{#2} }

% A three-argument union.
\newcommand*{\unionthree}[3]{ \union{\union{#1}{#2}}{#3} }

% The indexed union of many things.
\newcommand*{\unionmany}[3]{ \binopmany{\bigcup}{#1}{#2}{#3} }

% The intersection of its two arguments.
\newcommand*{\intersect}[2]{ {#1}\cap{#2} }

% A three-argument intersection.
\newcommand*{\intersectthree}[3]{ \intersect{\intersect{#1}{#2}}{#3} }

% The indexed intersection of many things.
\newcommand*{\intersectmany}[3]{ \binopmany{\bigcap}{#1}{#2}{#3} }

% The Cartesian product of two things.
\newcommand*{\cartprod}[2]{ {#1}\times{#2} }

% The Cartesian product of three things.
\newcommand*{\cartprodthree}[3]{ \cartprod{{#1}}{\cartprod{{#2}}{{#3}}} }

% The indexed Cartesian product of many things.
\newcommand*{\cartprodmany}[3]{ \binopmany{\bigtimes}{#1}{#2}{#3} }


\fi

%
% The standard corollary, definition, theorem, etc. that we all define
% in every single paper.
%
\ifx\havemjotheorem\undefined
\def\havemjotheorem{1}

\ifx\theoremstyle\undefined
  \usepackage{amsthm}
\fi

\theoremstyle{plain}
\newtheorem{corollary}{Corollary}
\newtheorem{exercise}{Exercise}
\newtheorem{lemma}{Lemma}
\newtheorem{proposition}{Proposition}
\newtheorem{theorem}{Theorem}

\theoremstyle{definition}
\newtheorem{definition}{Definition}
\newtheorem{example}{Example}

\theoremstyle{remark}
\newtheorem{remark}{Remark}


\fi


% Cleveref with Oxford comma!
\usepackage[capitalize,nameinlink,noabbrev,nosort]{cleveref}
\newcommand{\creflastconjunction}{, and\nobreakspace}

% Support for "inequality", "implication", "condition", and "form"
% references in addition to the usual "equation" references.
\crefname{inequality}{Inequality}{Inequalities}
\creflabelformat{inequality}{#2\textup{(#1)}#3}
\crefname{implication}{Implication}{Implications}
\creflabelformat{implication}{#2\textup{(#1)}#3}
\crefname{condition}{Condition}{Conditions}
\creflabelformat{condition}{#2\textup{(#1)}#3}
\crefname{form}{Form}{Forms}
\creflabelformat{form}{#2\textup{(#1)}#3}

% amsthm must be loaded BEFORE clevered, but the \newtheorem goes AFTER
\theoremstyle{plain}
\newtheorem{conjecture}{Conjecture}

% \Rnplus[1], but with a visible "1"
\newcommand*{\Roneplus}[0]{ \mathbb{R}^{1}_{+} }

% The signature of its argument.
\newcommand*{\signature}[1]{ \sigma\of{#1} }

\newcommand*{\smallestidx}[0]{\iota}

\begin{document}
  \title{
    The uniqueness of Lyapunov rank among symmetric cones
  }
  \author{Michael Orlitzky \and Giovanni Barbarino}
  %% \email{michael@orlitzky.com}
  %% \address{University of Maryland Baltimore County,
  %%          1000 Hilltop Circle, Baltimore, MD 21250}

  \maketitle

  \begin{abstract}
    The Lyapunov rank of a cone is the dimension of the Lie algebra of
    its automorphism group. It is invariant under linear isomorphism
    and in general not unique---two or more non-isomorphic cones can
    share the same Lyapunov rank. It is therefore not possible in
    general to identify cones using Lyapunov rank. But what if we look
    only among \emph{symmetric} cones? Are there any that can be
    uniquely identified (up to isomorphism) by their Lyapunov ranks?
    We provide a complete answer for irreducible cones and make some
    progress in the general case.
  \end{abstract}

  \begin{section}{Introduction}
    The motivation for Lyapunov rank comes from linear programming:
    the familiar complementary slackness condition $\ip{x}{s} = 0$ for
    $x,s \in \Rnplus$ can be separated into $n$ equations
    $\setc{x_{i}s_{i} = 0}{i = 1,2,\ldots,n}$ using the nonnegativity
    of the components. The idea is that these $n$ equations should be
    easier to solve than the one equation we started
    with. Complementary slackness is a necessary condition for
    optimality, so this helps us find candidate solutions. In this
    example, $n$ is the dimension of the space of all $n$-by-$n$
    diagonal matrices. It is not terribly obvious, but that space is
    the Lie algebra of $\Aut{\Rnplus}$.

    Albeit under a different name, the study of Lyapunov rank goes
    back to Rudolf, Noyan, Papp, and
    Alizadeh~\cite{rudolf-noyan-papp-alizadeh} who realized that this
    scenario generalizes to a cone program over a proper cone $K$ and
    its dual $\dual{K}$. In a cone program, complementary slackness is
    still a necessary condition for optimality, and the the condition
    $\ip{x}{s} = 0$ for $x \in K$ and $s \in \dual{K}$ can always be
    separated into a system of equations via ``bilinear
    complementarity relations.'' The number of equations, denoted
    $\lyapunovrank{K}$, depends only on the cone $K$ and is called the
    bilinearity rank (by them) or the Lyapunov rank (by us) of $K$. If
    you plan to search for candidate solutions by decomposing the
    complementary slackness condition, Lyapynov rank is a measure of
    how cooperative the cone will be.

    If the cone $K$ lives in an $n$-dimensional space, then the set of
    all candidate pairs $\setc{\pair{x}{s}}{x \in K, s \in \dual{K},
      \ip{x}{s} = 0}$ forms an $n$-dimensional
    manifold~\cite{rudolf-noyan-papp-alizadeh}. We might therefore
    hope to obtain $n$ or more equations ($\lyapunovrank{K} \ge n$) to
    be solved simultaneously. Rudolf, Noyan, Papp, and Alizadeh call
    such a cone $K$ \emph{bilinear}, and show that several important
    cones are not bilinear. A few years later, Gowda and Tao were able
    to connect these ideas to the Lyapunov-like operators on
    $K$~\cite{gowda-tao_-_rank}. Lyapunov-like operators on $K$ form
    the Lie algebra of the automorphism group of $K$, a vector space
    whose dimension is $\lyapunovrank{K}$. Linearly independent
    Lyapunov-like operators correspond to ``extra'' equations, and
    vice-versa. For this reason the name Lyapunov rank was chosen.

    Around the same time, Seeger and Sossa coined the term
    \emph{Loewnerian cone} for a cone that is isomorphic to the real
    symmetric PSD cone, $\Snplus$. As demonstrated by Seeger and
    Sossa, Loewnerian cones have important applications to
    complementarity
    problems~\cite{seeger-sossa_-_complementarity_problems}. But how
    do we know if a cone is Loewnerian? Sossa, Gowda, Tao, and one of
    the authors would eventually meet at the University of Maryland
    Baltimore County. The question naturally arose: can Lyapunov rank
    be used to show that a given cone is, or is not, isomorphic to
    $\Snplus$?  If the Lyapunov rank of $K$ differs from that of
    $\Snplus$, then $K$ and $\Snplus$ are not isomorphic. But perhaps
    there exist non-Loewnerian cones $K$ such that $\lyapunovrank{K} =
    \lyapunovrank{\Snplus}$?

    It turns out that we are able to answer the question for $\Snplus$
    negatively and with relative ease, but the problem inspires us to
    investigate further. Specifically, we ask whether or not there are
    any symmetric (self-dual and homogeneous) cones whose Lyapunov
    ranks are not shared by any other non-isomorphic symmetric
    cone. This question is not so easy to resolve, and the answer is
    somewhat unexpected.
  \end{section}

  \begin{section}{Background}
    Our main objects of interest are symmetric cones in finite
    dimensions. All such cones arise as the cone of squares in some
    Euclidean Jordan algebra~\cite{faraut-koranyi}. They are
    self-dual, and in particular, full-dimensional---they span the
    ambient space. This will be important because if we are searching
    for symmetric cones isomorphic to a given cone, it is a waste of
    time to check inside spaces of higher/lower dimensions.

    Every symmetric cone is the orthogonal direct sum of a unique set
    of nontrivial irreducible symmetric cones (Faraut and
    Korányi~\cite{faraut-koranyi}, Proposition~III.4.5). Moreover, up
    to isomorphism, there are only five families of irreducible
    factors. Chapter~V of Faraut and Korányi classifies the simple
    Euclidean Jordan algebras; the classification of irreducible
    symmetric cones then follows from the correspondence between
    symmetric cones and Euclidean Jordan algebras. Starting from a
    symmetric cone, one obtains a Euclidean Jordan algebra (Faraut and
    Korányi, Chapter~III) in which the given cone is the cone of
    squares. That algebra is the orthogonal direct sum of a unique set
    of nontrivial simple algebras by Proposition~III.4.4. In each of
    those simple factors, the cone of squares is a symmetric cone,
    which must be irreducible to avoid contradictions.

    \begin{theorem}\label{thm:irreducible families}
      % Take an empty sum to get K = {0}. This allows us to stipulate
      % that n >= 1 for the spin algebras and to say that the factors
      % are nontrivial.
      Every symmetric cone is the orthogonal direct sum of a unique
      set of nontrivial irreducible symmetric cones, and every
      nontrivial irreducible symmetric cone is Jordan-isomorphic to a
      member of one of the five families,
      %
      \begin{enumerate}
        \begin{item}
          The $n$-dimensional Lorentz cones, $\Lnplus$, for $n \ge 1$,
        \end{item}
        \begin{item}
          The $n \times n$ real symmetric positive semidefinite cones,
          $\Hnplus\of{\Rn[1]}$, for $n \ge 3$,
        \end{item}
        \begin{item}
          The $n \times n$ complex Hermitian positive semidefinite
          cones $\Hnplus\of{\Cn[1]}$, for $n \ge 3$,
        \end{item}
        \begin{item}
          The $n \times n$ quaternion Hermitian positive semidefinite
          cones\footnote{These cones of squares were identified as
          positive-semidefinite cones only
          recently~\cite{orlitzky_-_jordan_automorphisms}.},
          $\Hnplus\of{\quaternions}$, for $n \ge 3$,
        \end{item}
        \begin{item}
          The cone of squares in the Albert algebra,
          $\Hnplus[3]\of{\octonions}$.
        \end{item}
      \end{enumerate}
      %
    \end{theorem}

    The restrictions on $n$ ensure that the families are distinct up
    to isomorphism: the cones $\Hnplus[2]\of{\Rn[1]}$,
    $\Hnplus[2]\of{\Cn[1]}$, $\Hnplus[2]\of{\quaternions}$, and
    $\Hnplus[2]\of{\octonions}$ are symmetric, but they are isomorphic
    to Lorentz cones of appropriate sizes (Faraut and
    Korányi~\cite{faraut-koranyi}, Corollary~IV.1.5). The trivial cone
    in the trivial space, obtained as an empty sum, is also symmetric
    but is spiritually a member of the family $\Lnplus$ with $n=0$.

    The Lyapunov rank of a closed convex cone (and in particular, of a
    symmetric cone) is the dimension of the Lie algebra of its
    automorphism group. For our purposes, it is simply a number
    associated with the cone, and happens to be additive on a direct
    sum. Given that Cartesian products and direct sums are linearly
    isomorphic, this follows from Lemma~1 and Proposition~9 of Rudolf,
    Noyan, Papp, and Alizadeh~\cite{rudolf-noyan-papp-alizadeh}.

    \begin{proposition}\label{prop:lyapunov rank is additive on direct sums}
      If $K$ and $J$ are symmetric cones, then
      $\lyapunovrank{\directsum{K}{J}} = \lyapunovrank{K} +
      \lyapunovrank{J}$.
    \end{proposition}

    Taken together, we call the combination of dimension and Lyapunov
    rank the \emph{signature} of a cone.

    \begin{definition}[signature, similacra]
      Let $K$ be a closed convex cone. The \emph{signature} of $K$ is
      the pair $\signature{K} \coloneqq
      \pair{\dim\of{K}}{\lyapunovrank{K}}$. If $J$ is another closed
      convex cone, we say that $J$ is a \emph{similacrum} of $K$ if
      $\signature{K} = \signature{J}$ and if $K$ and $J$ are
      non-isomorphic. The plural of similacrum is \emph{similacra}. As
      shorthand we write $K \sim J$ to indicate that $K$ and $J$ are
      similacra of one another. If they are not, then instead we write
      $K \nsim J$.
    \end{definition}

    To be clear, by ``isomorphic,'' we mean the following: two closed
    convex cones $K \subseteq V$ and $J \subseteq W$ are isomorphic
    (and we write $K \cong J$) if there exists an invertible linear $A
    : V \to W$ such that $A\of{K} = J$. If two Euclidean Jordan
    algebras are Jordan-isomorphic (a stronger notion), then their
    respective cones of squares are necessarily isomorphic in the
    present sense.

    Gowda and Tao~\cite{gowda-tao_-_rank} computed Lyapunov ranks for
    the five irreducible families appearing in \cref{thm:irreducible
      families}. Using conjugate-symmetry and the fact that
    $\Cn[1],\quaternions,\octonions$ are $2,4,8$-dimensional algebras
    over $\Rn[1]$, it is easy to deduce their dimensions (and thus,
    their signatures) as well. We list these below. Combined with
    \cref{prop:lyapunov rank is additive on direct sums}, and waving
    away the practical difficulty of computing the direct sum
    decomposition in the first place, this gives the signature of any
    symmetric cone.

    \begin{table}[H]
      \begin{center}
        \begin{tabular}{c cc}
          $K$ & $\dim\of{K}$ & $\lyapunovrank{K}$\\
          \midrule
          $\Lnplus$                   & $n$
                                      & $\frac{n^{2} - n + 2}{2}$
          \\
          $\Hnplus\of{\Rn[1]}$        & $\frac{n^{2}+n}2$
                                      & $n^{2}$
          \\
          $\Hnplus\of{\Cn[1]}$        & $n^{2}$
                                      & $2n^{2} - 1$
          \\
          $\Hnplus\of{\quaternions}$  & $2n^{2} - n$
                                      & $4n^{2}$
          \\
          $\Hnplus[3]\of{\octonions}$ & $27$
                                      & $79$
        \end{tabular}
        \caption{Signatures of irreducible symmetric cones for $n \ge 2$}\label{tab:irreducible signatures}
      \end{center}
    \end{table}

    To this list we add one more, the nonnegative orthant in $n$
    dimensions:

    \begin{table}[H]
      \begin{center}
        \begin{tabular}{c cc}
          $K$ & $\dim\of{K}$ & $\lyapunovrank{K}$\\
          \midrule
          $\Rnplus$ & $n$ & $n$
        \end{tabular}
        \caption{Signature of $\Rnplus$ for $n \ge 0$}\label{tab:rnplus signature}
      \end{center}
    \end{table}

    \noindent
    The cone $\Rnplus = \directsummany{i=1}{n}{\Lnplus[1]}$ is
    symmetric but generally not irreducible (unless $n = 0$ or $n =
    1$). We include it because it is easier, for example, to write
    $\Rnplus[3]$ than it is to write $\Lnplus[1] \oplus \Lnplus[1]
    \oplus \Lnplus[1]$. In dimension $n \le 2$, all proper cones are
    polyhedral and therefore isomorphic to $\Rnplus$. As a result,
    when we are working up to isomorphism, we will usually write
    $\Roneplus$ or $\Rnplus[2]$ to indicate a one- or two-dimensional
    symmetric cone.
  \end{section}

  \begin{section}{Irreducible symmetric cones}\label{sec:irreducible cones}
    As a starting point, we show that $\Lnplus$ has no symmetric
    similacra, regardless of $n$. This is not too surprising; it
    follows from the fact that the Lyapunov rank of $\Lnplus$ is
    strictly maximal in $n$ dimensions. The latter is known to most
    people working with Lyapunov rank, and the details can be worked
    out using \cref{tab:irreducible signatures}.

    \begin{proposition}\label{prop:rank of lorentz is maximal}
      Among non-isomorphic $n$-dimensional symmetric cones, the
      Lyapunov rank of $\Lnplus$ is strictly maximal for all $n \ge
      0$.
      %
      \begin{proof}
        First we note that if $n \le 2$, then all symmetric cones have
        exactly $n$ extreme rays and are therefore isomorphic to
        $\Rnplus$. The statement is therefore trivially true for $n
        \le 2$. We next show that splitting $\Lnplus$ into two smaller
        Lorentz cones will necessarily reduce the Lyapunov
        rank. Taking any $k \in \set{1,2,\ldots,n-1}$, and assuming
        now that $n \ge 3$, we can simply compute:
        %
        \begin{equation*}
          \lyapunovrank{\Lnplus}
          -
          \lyapunovrank{\directsum{\Lnplus[k]}{\Lnplus[n-k]} }
          =
          \qty{n-k}{k} - 1.
        \end{equation*}
        %
        For $n \ge 3$, this difference is strictly positive, showing
        that the direct sum has a strictly inferior Lyapunov rank.

        Finally, we show that members of the other four irreducible
        families have Lyapunov ranks smaller than a Lorentz cone of
        the same dimension. There's only one cone of octonion matrices
        to worry about, and it's easy to see that
        $\lyapunovrank{\Lnplus[27]} >
        \lyapunovrank{\Hnplus[3]\of{\octonions}}$. For the other three
        families, we compute:
        %
        \begin{alignat*}{3}
          \lyapunovrank{\Lnplus[\qty{m^{2}+m}/2]}
          &-&
          \lyapunovrank{\Hnplus[m]\of{\Rn[1]}}
          &=
          \frac{1}{8}m^{4} + \frac{1}{4}m^{3} - \frac{9}{8}m^{2} - \frac{1}{4}m + 1,
          \\
          \lyapunovrank{\Lnplus[m^{2}]}
          &-&
          \lyapunovrank{\Hnplus[m]\of{\Cn[1]}}
          &=
          \frac{1}{2}m^{4} - \frac{5}{2}m^{2} + 2,
          \\
          \lyapunovrank{\Lnplus[2m^{2} - m]}
          &-&
          \lyapunovrank{\Hnplus[m]\of{\quaternions}}
          &=
          2m^{4} - 2m^{3} - \frac{9}{2}m^{2} + \frac{1}{2}m + 1.
        \end{alignat*}
        %
        At $m = 3$, the higher-order (positive) terms already dominate
        the lower-order (negative) ones. These expressions therefore
        remain strictly positive as $m$ grows. The result now follows:
        in $n \ge 3$, splitting the Lorentz cone into two smaller
        Lorentz cones will decrease the Lyapunov rank, and using
        anything other than Lorentz cones will only decrease it
        further.
      \end{proof}
    \end{proposition}

    \begin{corollary}\label{cor:lorentz cone has no similacra}
      $\Lnplus$ has no symmetric similacra for any $n \ge 0$.
    \end{corollary}

    Looking beyond symmetric cones for a moment, it is known that the
    Lyapunov rank of the Lorentz cone is maximal among all proper
    cones of the same dimension~\cite{orlitzky_-_tight_bounds}, but it
    is not known if the inequality is strict. In other words, whether
    or not there is a non-isomorphic, $n$-dimensional, proper (but not
    symmetric) cone having the same Lyapunov rank as $\Lnplus$ is
    still open.

    \begin{proposition}\label{prop:real psd similacra}
      $\Hnplus\of{\Rn[1]}$ has symmetric similacra for all $n \ge 3$.
      %
      \begin{proof}
        Compare the signature of $\Lnplus[n+1] \oplus
        \Rnplus[\qty{n^{2} - n - 2}/2]$.
      \end{proof}
    \end{proposition}

    Interestingly, the formula in \cref{prop:real psd similacra}
    produces a duplicate signature for $n=2$ as well. The resulting
    cone $\Lnplus[3]$ is however isomorphic to $\Hnplus[2]\of{\Rn[1]}$
    and therefore not a similacrum.

    \begin{proposition}\label{prop:uniqueness of h3}
      $\Hnplus[3]\of{\Cn[1]}$ has no symmetric similacra.

      \begin{proof}
        Keeping in mind that $n \ge 3$ is required to produce
        non-isomorphic families, the dimension of this cone rules out
        $\Hnplus\of{\Cn[1]}$, $\Hnplus\of{\quaternions}$ and
        $\Hnplus[3]\of{\octonions}$ as factors. By \cref{prop:real psd
          similacra} we may therefore assume without loss of
        generality that all factors are Lorentz cones of various
        sizes. Since $\lyapunovrank{\Lnplus} > 17$ for $n > 6$, we
        need only consider Lorentz cone factors of dimension six or
        less. And $\lyapunovrank{\Lnplus[6]} = 16$, so if $\Lnplus[6]$
        were one factor, then the other factors would have combined
        signature $\pair{3}{1}$, which is not possible. So, all
        Lorentz cone factors have dimension five or less.

        If $\Lnplus[5]$ is a factor, then the remaining factors
        (combined) have signature $\pair{4}{6}$. The Lyapunov ranks of
        $\Rnplus[4]$ and $\Lnplus[3] \oplus \Roneplus$
        are too small, and that of $\Lnplus[4]$ is too large. Thus, a
        similacrum cannot have $\Lnplus[5]$ factors.

        If $\Lnplus[4]$ is a factor, then the other factors have
        signature $\pair{5}{10}$. $\Lnplus[5]$ has already been ruled
        out as a factor, and the Lyapunov rank of $\Lnplus[4] \oplus
        \Roneplus$ is too small. \cref{prop:rank of lorentz is % chktex 12
          maximal} shows that further splitting of $\Lnplus[4]$ will
        not help, so $\Lnplus[4]$ cannot be a factor.

        Finally, we note that $\lyapunovrank{\Lnplus[3] \oplus
          \Lnplus[3] \oplus \Lnplus[3]}$ is too small.
      \end{proof}
    \end{proposition}

    \begin{proposition}\label{prop:complex PSD of size four or more has similacra}
      $\Hnplus\of{\Cn[1]}$ has symmetric similacra for all $n \ge 4$.

      \begin{proof}
        If $n = 4$, then $\Hnplus[4]\of{\Cn[1]} \sim \Lnplus[5] \oplus
        \Lnplus[5] \oplus \Lnplus[4] \oplus \Rnplus[2]$. And if $n \ge 5$,
        behold:
        %
        \begin{equation*}
          \Hnplus\of{\Cn[1]}
          \sim \Lnplus[n+1]
               \oplus \Lnplus[n+1]
               \oplus \Rnplus[\qty{n^{2} - 5n + 1}]
               \oplus \underbrace{\Lnplus[3] \oplus \cdots \oplus \Lnplus[3]}_{n-1 \text{ times}}.
          \qedhere
        \end{equation*}
        %
      \end{proof}
    \end{proposition}

    \begin{proposition}\label{prop:quaternion psd has similacra}
      $\Hnplus\of{\quaternions}$ has symmetric similacra for all $n \ge 3$.

      \begin{proof}
        The first few examples we find by trial and error:
        %
        \begin{align*}
          \Hnplus[3]\of{\quaternions}
            &\sim
            \directsum{\Lnplus[8]}{\Rnplus[7]},\\
          \Hnplus[4]\of{\quaternions}
            &\sim
            \Lnplus[10] \oplus \Rnplus[18],\\
            \Hnplus[5]\of{\quaternions}
            &\sim
            % The weirdest fucking thing happens if you un-brace the
            % "3" in "33".
            \directsum{\Lnplus[12]}{\Rnplus[{3}3]}.
        \end{align*}
        %
        For $n \ge 6$, witness:
        %
        \begin{equation*}
          \Hnplus[n]\of{\quaternions}
          \sim
          \Hnplus[n+1]\of{\Cn[1]}
               \oplus \Lnplus[n+1]
               \oplus \Lnplus[n+1]
               \oplus \Rnplus[\qty{n^{2} - 5n - 3}].
               \qedhere
        \end{equation*}
        %
      \end{proof}
    \end{proposition}


    \begin{proposition}\label{prop:octonion psd has similacra}
      $\Hnplus[3]\of{\octonions}$ has symmetric similacra.

      \begin{proof}
        Compare the signature of $\Lnplus[11] \oplus \Lnplus[5] \oplus
        \Lnplus[3] \oplus \Rnplus[8]$.
      \end{proof}
    \end{proposition}

    Combining all of these, we conclude the section with our first
    major result.

    \begin{theorem}\label{thm:only irreducible factors}
      If $K$ is an irreducible symmetric cone having no symmetric
      similacra, then either $K = \Hnplus[3]\of{\Cn[1]}$ or $K =
      \Lnplus$ for some $n \ge 0$.
    \end{theorem}

    The similacra in this section appear to have been pulled from thin
    air. Which is not wholly untrue, but with the benefit of
    hindsight, we can explain how one might have found them. Searching
    for similacra is a recursive problem. If $K$ is the target cone,
    then we are looking to achieve $\signature{K} =
    \pair{\dim\of{K}}{\lyapunovrank{K}}$. It may not be obvious how to
    do that, but one thing we can try is to suppose (for example) that
    $\Hnplus\of{\Rn[1]}$ is one of the factors. That factor will
    contribute $\qty{n^{2} + n}/2$ to the first component of the
    signature and $n^{2}$ to the second, meaning that now the problem
    is reduced to finding a signature of $\pair{\dim\of{K} -
      \qty{n^{2} + n}/2}{\lyapunovrank{K} - n^{2}}$. If we're lucky,
    this problem will be easier to solve. If not, we can go back,
    assume that some other factor is present in the similacrum, and
    try again.

    To bring that process to a close, it helps to have collection of
    solutions to subproblems---means of achieving a certain family of
    signatures, like $\pair{3n}{4n}$ with $n \in \Nn[1]$. If you can
    construct a cone whose dimension is $3n$ and whose Lyapunov rank
    is $4n$, then you are done if the signature $\pair{3n}{4n}$ arises
    as a subproblem. It turns out that, depending on how you count
    them, only two or three subproblems are needed to construct the
    similacra in this section.

    The first is straightforward: in \cref{prop:complex PSD of size
      four or more has similacra}, we used the fact that
    $\signature{\Lnplus[3]} = \pair{3}{4}$ which leads a solution of
    the $\pair{3n}{4n}$-subproblem,
    %
    \begin{equation*}
      \signature{
        \underbrace{\Lnplus[3] \oplus \cdots \oplus \Lnplus[3]}_{n \text{ times}}
      }
      =
      \sum_{i=1}^{n}
      \signature{\Lnplus[3]}
      =
      \pair{3n}{4n}.
    \end{equation*}
    %
    For the second, fix a dimension $d$, and let $K =
    \directsum{\Lnplus[m]}{\Rnplus[d - m]}$, so that $\dim\of{K} =
    d$. The Lyapunov rank of $K$ is,
    %
    \begin{align*}
      \lyapunovrank{K}
        &= \frac{m^{2} - m + 2}{2} + \qty{d - m}\\
        &= d + \frac{m^{2} - 3m + 2}{2}\\
        &= d + \frac{1}{2}\qty{m-1}\qty{m-2}.
    \end{align*}
    %
    Thus we have a solution for the $\pair{d}{d +
      \frac{1}{2}\qty{m-1}\qty{m-2}}$-subproblem, for all $d$, and for
    all $m \le d$. Moreover, if $d - m \ge m$, we can repeat this
    process with another $\Lnplus[m]$ factor to eliminate the $1/2$:
    %
    \begin{equation*}
      \lyapunovrank{\Lnplus[m] \oplus \Lnplus[m] \oplus \Rnplus[d - 2m]}
        = d + \qty{m-1}\qty{m-2}.
    \end{equation*}
    %
    This lets us solve the $\pair{d}{d +
      \qty{m-1}\qty{m-2}}$-subproblem as well, regardless of $d$, and
    assuming only that $2m \le d$. To see this in action, consider
    $\Hnplus\of{\Rn[1]}$. Its dimension is $\qty{n^{2} + n}/2$ and its
    Lyapunov rank is $n^{2}$, so its Lyapunov rank is $\qty{n^{2} -
      n}/2$ greater than its dimension. If we want to apply the
    subproblem solution, we need to find an $m$ such that
    $\qty{m-1}\qty{m-2}/2 = \qty{n^{2} - n}/2$. But this is easy: take
    $m \coloneqq n + 1$. The solution to the subproblem then gives the
    similacra in \cref{prop:real psd similacra}. The second form of
    this subproblem is used in \cref{prop:quaternion psd has
      similacra} after guessing that one factor is
    $\Hnplus[n+1]\of{\Cn[1]}$.
  \end{section}

  \begin{section}{Reducible cones}\label{sec:reducible cones}
    To venture beyond irreducible cones makes our job much more
    difficult. We already have several examples of non-irreducible
    cones whose Lyapunov ranks are not unique. Every similacrum we've
    constructed so far has symmetric similacra---that's why we
    constructed it! The best we can hope for is to narrow down the
    possibilities a bit.

  \begin{lemma}\label{lem:no multiple complex factors}
      If a symmetric cone contains more than one
      $\Hnplus[3]\of{\Cn[1]}$ factor (up to isomorphism), then it has
      symmetric similacra.
      %
      \begin{proof}
        Let,
        \begin{align*}
          K_{2} &\coloneqq \Lnplus[7]
                 \oplus    \Lnplus[3]
                 \oplus    \Rnplus[8],\\
          K_{3} &\coloneqq \Lnplus[8]
                 \oplus    \Lnplus[4]
                 \oplus    \Rnplus[15].
        \end{align*}
        %
        Then,
        %
        \begin{align*}
          \signature{K_{2}}
            = \pair{18}{34}
            &= 2\signature{\Hnplus[3]\of{\Cn[1]}},\\
          \signature{K_{3}}
            = \pair{27}{51}
            &= 3\signature{\Hnplus[3]\of{\Cn[1]}}.
        \end{align*}
        %
        Suppose our target cone contains $m$ copies of
        $\Hnplus[3]\of{\Cn[1]}$. Since $m = 2k + 3j$ for some $k \in
        \Nn[1]$ and $j \in \set{0,1}$, we can construct a symmetric cone
        containing $k$ copies of $K_{2}$ and $j$ copies of $K_{3}$ (in
        lieu of the $m$ copies of $\Hnplus[3]\of{\Cn[1]}$) sharing its
        signature. The new cone cannot be isomorphic to the original
        because the direct sum decomposition into a set of irreducible
        factors is unique up to isomorphism.
      \end{proof}
    \end{lemma}

    \begin{theorem}\label{thm:two unique families}
      Suppose $K$ is a symmetric cone. If $K$ has no symmetric
      similacra, then $K$ is isomorphic to
      %
      \begin{equation*}
        \qty{\Hnplus[3]\of{\Cn[1]}}^{j}
        \oplus
        \Lnplus[n_{1}]
        \oplus
        \Lnplus[n_{2}]
        \oplus
        \cdots
        \oplus
        \Lnplus[n_{k}],
      \end{equation*}
      %
      where $j \in \set{0,1}$ and $k, n_{i} \in \Nn[1]$. If, on the
      other hand, $K$ does have symmetric similacra, then it has a
      similacrum of the above form.

      \begin{proof}
        \cref{thm:only irreducible factors} shows that the other
        irreducible families have similacra involving
        $\Hnplus[3]\of{\Cn[1]}$ and $\Lnplus[n_{i}]$, and \cref{lem:no
          multiple complex factors} shows that we need (or can have)
        at most one $\Hnplus[3]\of{\Cn[1]}$ factor.
      \end{proof}
    \end{theorem}

    This result can be interpreted as saying that every symmetric cone
    shares its signature with a direct sum of Lorentz cones and/or
    $\Hnplus[3]\of{\Cn[1]}$ via either isomorphism or similacra. This
    much is true, but we caution that even this ``reduced''
    representation of a similacrum is not unique.

    \begin{example}\label{ex:reduction is not unique}
      It is easy to check that
      %
      \begin{equation*}
        \Lnplus[4] \oplus \Lnplus[3] \oplus \Lnplus[3] \oplus \Lnplus[3]
        \sim
        \Lnplus[5] \oplus \Rnplus[8].
      \end{equation*}
      %
      As a result, we can derive a second (non-ismomorhpic) similacrum
      consisting only of Lorentz cones in \cref{prop:octonion psd has
        similacra}.
    \end{example}

    Using \cref{thm:two unique families} we can reduce the search for
    similacra to essentially two integer partition problems, one for
    $j = 0$ and one for $j = 1$ in \cref{thm:two unique
      families}. Suppose $\signature{K} =
    \pair{\dim\of{K}}{\lyapunovrank{K}}$ is given. One possibility is
    that $K$ has a similacrum with no $\Hnplus[3]\of{\Cn[1]}$ factor;
    that is, that
    %
    \begin{equation*}
      K
      \sim
      \Lnplus[n_{1}]
        \oplus
        \Lnplus[n_{2}]
        \oplus
        \cdots
        \oplus
        \Lnplus[n_{k}]
    \end{equation*}
    %
    where $\set{n_{i}}_{i=1}^{k}$ is an integer partition of
    $\dim\of{K}$ and
    %
    \begin{equation*}
      \lyapunovrank{K}
      =
      \sum_{i=1}^{k} \lyapunovrank{\Lnplus[n_{i}]}
      =
      \sum_{i=1}^{k} \qty{n_{i}^{2} - n_{i} + 2}/2.
    \end{equation*}
    %
    The set of all such integer partitions can be enumerated, and the
    corresponding sums are easily checked against
    $\lyapunovrank{K}$. The other option is that $K$ has a similacrum
    with exactly one $\Hnplus[3]\of{\Cn[1]}$ factor. In that case, we
    can subtract $\signature{\Hnplus[3]\of{\Cn[1]}} = \pair{9}{17}$ from
    $\signature{K}$ and repeat the procedure above. All \emph{other}
    factors will be Lorentz cone factors, so the problem is reduced to
    finding an integer partition $\set{n_{i}}_{i=1}^{k}$ of $\dim\of{K}
    - 9$ such that $\lyapunovrank{K} - 17 = \sum_{i=1}^{k}
    \qty{n_{i}^{2} - n_{i} + 2}/2$. This leads to a straightforward
    algorithm for checking examples of reasonably small dimension using
    (say) the accelerated integer partitioning scheme of Kelleher and
    O'Sullivan~\cite{kelleher-osullivan}.

    \cref{lem:no multiple complex factors} shows that if a symmetric
    cone has one $\Hnplus[3]\of{\Cn[1]}$ factor, then it cannot have any
    additional $\Hnplus[3]\of{\Cn[1]}$ factors without acquiring
    symmetric similacra. Does something similar happen when we combine
    $\Hnplus[3]\of{\Cn[1]}$ and $\Lnplus$ factors? It is worth
    investigating due to the following (rather obvious) result.

    \begin{lemma}\label{lem:subcone similacra}
      If any symmetric subcone of a symmetric cone $K$ has symmetric
      similacra, then so does $K$.
    \end{lemma}

    If we can determine which pairs of irreducible factors induce
    symmetric similacra, then we can rule out those pairs as factors in
    any larger symmetric cone that does not have symmetric
    similacra.

    \begin{lemma}\label{lem:no bigger factors}
      If $n \ge 10$, then $\directsum{\Hnplus[3]\of{\Cn[1]}}{\Lnplus}
      \nsim \directsum{J}{\Lnplus[n+k]}$ for any symmetric $J$ and any
      $k \ge 0$.
      %
      \begin{proof}
        We can rule out $k = 0$, because in that case, we are left with
        just $\Hnplus[3]\of{\Cn[1]}$ which has no similacra. With $k \ge
        1$, the smallest possible Lyapunov rank of any similacrum is at
        $k = 1$ with $J = \Rnplus[8]$. But if $n \ge 10$, then
        %
        \begin{equation*}
          \lyapunovrank{\Rnplus[8] \oplus \Lnplus[n+1]}
          -
          \lyapunovrank{\Hnplus[3]\of{\Cn[1]} \oplus \Lnplus}
          =
          n-9
          \ge
          1.
        \end{equation*}
        %
        Thus, the smallest possible Lyapunov rank is too large to work.
      \end{proof}
    \end{lemma}

    \begin{proposition}\label{prop:no similacra for h3plus and big enough lnplus}
      If $n \ge 31$, then $\directsum{\Hnplus[3]\of{\Cn[1]}}{\Lnplus}$
      has no symmetric similacra.

      \begin{proof}
        We first note that it suffices to search for a similacrum of the
        form $\directsummany{i=1}{k}{\Lnplus[n_{i}]}$ with $n_{1} \ge
        n_{2} \ge \cdots \ge n_{k}$ and $k \ge 2$. If a similacrum
        contains $\Hnplus[3]\of{\Cn[1]}$, then, by subtracting
        $\pair{9}{17}$ from both sides, it also contains a similacrum
        for $\Lnplus$, which is not possible by \cref{cor:lorentz cone
          has no similacra}. The same corollary makes $k = 1$
        impossible, and we may reorder components to put the
        superscripts in whatever order we like. In light of \cref{lem:no
          bigger factors}, we may further assume that $n_{1} \le n-1$,
        and thus $n_{i} \le n-1$ for all $i \in \set{1,2,\ldots,k}$.

        The problem now is to find an integer partition
        $\set{n_{1},n_{2},\ldots,n_{k}}$ of $n+9$ such that
        $\sum_{i=1}^{k}\lyapunovrank{\Lnplus[n_{i}]} =
        \lyapunovrank{\Lnplus} + 17$. To simplify the notation, we let
        %
        \begin{equation*}
          f\of{x} \coloneqq \lyapunovrank{\Lnplus[x]},
        \end{equation*}
        %
        and note that, by \cref{prop:rank of lorentz is maximal},
        %
        \begin{equation}\label[inequality]{ineq:f is superadditive}
          f\of{x} + f\of{y} \le f\of{x + y}.
        \end{equation}
        %
        It is also easy to see that
        %
        \begin{equation}\label{eq:f of n - 1}
          f\of{n - 1} = f\of{n} - \qty{n - 1}.
        \end{equation}
        %
        If we fix $\delta > 0$, then $g_{\delta}\of{x} \coloneqq f\of{x}
        - f\of{x - \delta}$ is increasing since $g_{\delta}'\of{x} =
        \delta$. In particular, $g_{\delta}\of{n - 1} \ge
        g_{\delta}\of{10+\delta}$ whenever $n - 1 \ge 10 + \delta$. In
        other words,
        %
        \begin{equation}\label[implication]{impl:growth inequality for f}
          \begin{gathered}
            n - 1 \ge 10 + \delta\\
            \implies\\
            f\of{n - 1} + f\of{10} \ge f\of{n - 1 - \delta} + f\of{10 + \delta}.
          \end{gathered}
        \end{equation}
        %
        Suppose we have a partition $\set{n_{1},n_{2},\ldots,n_{k}}$ of
        $n+9$. We split the proof into two cases based on the dimension
        of the largest factor, $n_{1}$. In each case we will see that
        $\sum_{i=1}^{k}f\of{n_{i}}$ cannot be equal to $f\of{n} + 17$.
        %
        \begin{pcases}
          \begin{case}[$n_{1} \le 9$]
            Let $S\of{j} \coloneqq \sum_{i=1}^{j}n_{i}$ with $S\of{0}
            \coloneqq 0$. By assumption and definition we may
            immediately conclude that $S\of{j} + 9 \ge S\of{j+1} >
            S\of{j}$ for all $j<k$, and that $S\of{k} = n + 9 \ge 40$.
            Define $\smallestidx$ to be the smallest index for which
            $S\of{\smallestidx} > 9$ holds. There is such an index,
            and $\smallestidx \ge 2$. It follows that $S\of{\smallestidx
              - 1} \le 9$, $S\of{\smallestidx} \le 18$, and $\smallestidx <
            k$. Applying \cref{ineq:f is superadditive} repeatedly, we
            find that
            %
            \begin{equation*}
              \sum_{i=1}^{k}f\of{n_{i}}
              =
              \sum_{i=1}^{\smallestidx}f\of{n_{i}}
              +
              \sum_{i=\smallestidx+1}^{k}f\of{n_{i}}
              \le
              f\of{S\of{\smallestidx}} + f\of{n + 9 - S\of{\smallestidx}}.
            \end{equation*}
            %
            Let $\delta \coloneqq S\of{\smallestidx} - 10$. Then
            $\delta \ge 0$ by definition, and $10 + \delta =
            S\of{\smallestidx} \le 18 \le n - 1$.  As a result, we may
            rewrite the expression above in terms of $\delta$, and
            apply \cref{impl:growth inequality for f} to conclude that
            %
            \begin{equation*}
              \sum_{i=1}^{k}f\of{n_{i}}
              \le
              f\of{10 + \delta} + f\of{n - 1 - \delta}
              \le
              f\of{n - 1} + f\of{10}.
            \end{equation*}
            %
            Finally, using \cref{eq:f of n - 1}, we arrive at
            %
            \begin{equation*}
              \sum_{i=1}^{k}f\of{n_{i}}
              \le
              f\of{n} - \qty{n-1} + f\of{10}
              \le
              f\of{n} - 30 + 46
              <
              f\of{n} + 17.
            \end{equation*}
          \end{case}
          \begin{case}[$n_{1} \ge 10$]
            Recall that $n_{1} \le n-1$, and let $\delta \coloneqq n - 1 -
            n_{1} \ge 0$. Then $n_{1} = n - 1 - \delta$ and $\sum_{i=2}^{k}
            n_{i} = 10 + \delta$, so using \cref{ineq:f is superadditive},
            %
            \begin{equation*}
              \sum_{i=1}^{k} f\of{n_{i}}
              \le
              f\of{n - 1 - \delta} + f\of{10 + \delta}.
            \end{equation*}
            %
            We have $n_{1} \ge 10$ by assumption, so from \cref{impl:growth
              inequality for f} we conclude that
            %
            \begin{equation*}
              f\of{n - 1 - \delta} + f\of{10 + \delta}
              \le
              f\of{n-1} + f\of{10}.
            \end{equation*}
            %
            Finally, \cref{eq:f of n - 1} and the assumption that $n \ge
            31$ now quickly lead to
            %
            \begin{equation*}
              \sum_{i=1}^{k} f\of{n_{i}}
              \le
              f\of{n} - n + 47
              <
              f\of{n} + 17. \qedhere
            \end{equation*}
            %
          \end{case}
        \end{pcases}
      \end{proof}
    \end{proposition}

    To catalogue the similacra of
    $\directsum{\Hnplus[3]\of{\Cn[1]}}{\Lnplus}$, it now suffices to
    check only the first thirty values of $n$. Using the integer
    partition approach we have calculated the cases $n = 1,2\ldots,30$
    and have listed their similacra below. Absent rows indicate the
    absence of similacra.

    \begin{equation*}
      \begin{array}{c ccc}
        K & \dim\of{K} & \lyapunovrank{K} & \sim\\
        \midrule
        \Hnplus[3]\of{\Cn[1]} \oplus \Lnplus[2] &
          11 &
          19 &
          \Lnplus[5] \oplus \Lnplus[3] \oplus \Lnplus[3]\\
        %
        \Hnplus[3]\of{\Cn[1]} \oplus \Lnplus[3] &
          12 &
          21 &
          \Lnplus[4] \oplus \Lnplus[4] \oplus \Lnplus[4]\\
        %
        \Hnplus[3]\of{\Cn[1]} \oplus \Lnplus[4] &
          13 &
          24 &
          \Lnplus[6] \oplus \Lnplus[3] \oplus \Rnplus[4]\\
        %
        \Hnplus[3]\of{\Cn[1]} \oplus \Lnplus[5] &
          14 &
          28 &
          \Lnplus[6] \oplus \Lnplus[4] \oplus \Lnplus[3] \oplus \Roneplus\\
        %
        \Hnplus[3]\of{\Cn[1]} \oplus \Lnplus[6] &
          15 &
          33 &
          \Lnplus[5] \oplus \Lnplus[5] \oplus \Lnplus[5]\\
        %
        \Hnplus[3]\of{\Cn[1]} \oplus \Lnplus[7] &
          16 &
          39 &
          \Lnplus[6] \oplus \Lnplus[6] \oplus \Lnplus[4]\\
        %
        \Hnplus[3]\of{\Cn[1]} \oplus \Lnplus[8] &
          17 &
          46 &
          \Lnplus[9] \oplus \Lnplus[3] \oplus \Rnplus[5]\\
        %
        \Hnplus[3]\of{\Cn[1]} \oplus \Lnplus[9] &
          18 &
          54 &
          \Lnplus[10] \oplus \Rnplus[8]\\
        %
        \Hnplus[3]\of{\Cn[1]} \oplus \Lnplus[10] &
          19 &
          63 &
          \Lnplus[9] \oplus \Lnplus[7] \oplus \Lnplus[3]\\
        %
        \Hnplus[3]\of{\Cn[1]} \oplus \Lnplus[15] &
          24 &
          123 &
          \Lnplus[14] \oplus \Lnplus[8] \oplus \Rnplus[2]\\
        %
        \Hnplus[3]\of{\Cn[1]} \oplus \Lnplus[18] &
          27 &
          171 &
          \Lnplus[14] \oplus \Lnplus[13]\\
        %
        \Hnplus[3]\of{\Cn[1]} \oplus \Lnplus[21] &
          30 &
          228 &
          \Lnplus[19] \oplus \Lnplus[11]\\
        %
        \Hnplus[3]\of{\Cn[1]} \oplus \Lnplus[22] &
          31 &
          249 &
          \Lnplus[21] \oplus \Lnplus[9] \oplus \Roneplus\\
        %
        \Hnplus[3]\of{\Cn[1]} \oplus \Lnplus[30] &
          39 &
          453 &
          \Lnplus[29] \oplus \Lnplus[10]\\
      \end{array}
    \end{equation*}

    \begin{proposition}\label{prop:h3plus plus lnplus characterization}
      $\Hnplus[3]\of{\Cn[1]} \oplus \Lnplus$ has symmetric similacra if
      and only if $n \in \set{2,3,\ldots,10,15,18,21,22,30}$.
    \end{proposition}

    \begin{proposition}
      If $K$ is a symmetric cone with an $\Hnplus[3]\of{\Cn[1]}$ factor
      and no symmetric similacra, then
      %
      \begin{equation*}
        K
        \cong
        \Hnplus[3]\of{\Cn[1]}
        \oplus
        \Lnplus[n_{1}]
        \oplus
        \Lnplus[n_{2}]
        \oplus
        \cdots
        \oplus
        \Lnplus[n_{k}],
      \end{equation*}
      %
      where $k \in \Nn[1]$ and $n_{i} \in \Nn[1]
      \setminus \set{2,3,\ldots,10,15,18,21,22,30}$.
      %
      \begin{proof}
        With \cref{lem:subcone similacra} in mind, start from
        \cref{thm:two unique families} and remove the known similacra
        from the table.
      \end{proof}
    \end{proposition}

    The meta-problem we are up against with reducible cones is that,
    since there are so many possibilities, it is not terribly
    interesting to start crossing them off the list one at a time. It
    was worth looking at $\Hnplus[3]\of{\Cn[1]}$ because it is
    ``unique.'' But where can we go from here?

    Our last result aims to eliminate an entire family of
    possibilities by reducing them to subproblems. The idea is that,
    if a given cone has an $\Lnplus$ factor for large enough $n$, then
    any similacrum must have that same factor. Those two factors can
    then be subtracted, leaving the subproblem for whatever is
    left. This is essentially a generalization of \cref{prop:h3plus
      plus lnplus characterization}, except that in \cref{prop:h3plus
      plus lnplus characterization} we are able to able to check the
    cases $n < 31$ explicitly, and we know that the subproblem for
    $\Hnplus[3]\of{\Cn[1]}$ has no solution.

    If the problem is to find similacra of $\directsum{\Lnplus}{K}$,
    we will impose three conditions on $n$ relative to $K$:

    \begin{enumerate}
      \begin{item}\label[condition]{itm:condition1}
        $n \ge 2 + \lyapunovrank{K} - \dim\of{K}$,
      \end{item}
      \begin{item}\label[condition]{itm:condition2}
        $n \ge 2 + \lyapunovrank{\Lnplus[1 + \dim\of{K}]} - \lyapunovrank{K}$,
      \end{item}
      \begin{item}\label[condition]{itm:condition3}
        $n \ge 10$.
      \end{item}
    \end{enumerate}

    Taken together, these conditions imply two others that we'll need.

    \begin{lemma}\label{lem:more conditions}
      If $K$ is a nontrivial symmetric cone and if $n \in \Nn[1]$
      satisfies \cref{itm:condition1,itm:condition2,itm:condition3},
      then $n \ge 2\dim\of{K}$ and $n > \dim\of{K}$.

      \begin{proof}
        By averaging the first two conditions and using the third to
        rule out small $n$, we derive the additional bound $n -
        2\dim\of{K} \ge 0$. And since $K$ is nontrivial, $2\dim\of{K}
        > \dim\of{K}$.
      \end{proof}
    \end{lemma}

    We next use \cref{itm:condition1} to rule out the possibility that
    $\directsum{\Lnplus}{K}$ shares its signature with some other cone
    having a larger Lorentz cone factor than $\Lnplus$.

    \begin{lemma}\label{lem:no bigger factors pt2}
      If $K$ is symmetric and if $n \ge 2 + \lyapunovrank{K} -
      \dim\of{K}$, then $ \signature{\directsum{J}{\Lnplus[n+k]}} \ne
      \signature{\directsum{K}{\Lnplus}}$ for any symmetric $J$ and
      any $k \ge 1$.
      %
      \begin{proof}
        With $k \ge 1$, the smallest possible Lyapunov rank achievable
        by $\directsum{J}{\Lnplus[n+k]}$ in dimension $\dim\of{K}+n$
        is at $k = 1$ with $J = \Rnplus[\dim\of{K} - k]$. But if $n
        \ge 2 + \lyapunovrank{K} - \dim\of{K}$, then
        %
        \begin{equation*}
          \lyapunovrank{\Rnplus[\dim\of{K} - 1] \oplus \Lnplus[n+1]}
          -
          \lyapunovrank{K \oplus \Lnplus}
          =
          \dim\of{K} - 1 - \lyapunovrank{K} + n
          \ge
          1.
        \end{equation*}
        %
        Thus, the smallest possible Lyapunov rank is too large to work.
      \end{proof}
    \end{lemma}

    The other conditions can now be used to rule out cones having only
    Lorentz cone factors, where each of those factors is smaller than
    $\Lnplus$.

    \begin{lemma}\label{lem:not all smaller factors}
      If $K$ is a symmetric cone and if $n \in \Nn[1]$ satisfies
      \cref{itm:condition1,itm:condition2,itm:condition3}, then
      %
      \begin{equation*}
        \lyapunovrank{\directsummany{i=1}{k}{\Lnplus[n_{i}]}}
        <
        \lyapunovrank{\directsum{\Lnplus}{K}}
      \end{equation*}
      %
      for any $k \in \Nn[1]$ and any $n_{1},n_{2},\ldots,n_{k} < n$
      such that $\sum_{i=1}^{k} n_{i} = \dim\of{K} + n$.

      \begin{proof}
        The reasoning here follows closely that of \cref{prop:no
          similacra for h3plus and big enough lnplus}, with the new
        bounds on $n$ supplied by
        \cref{itm:condition1,itm:condition2,itm:condition3}. We
        quickly sketch the idea. As before, by rearranging factors, we
        may assume that $n_{1} \ge n_{2} \ge \cdots \ge n_{k}$.

        Recall the definition $f\of{x} \coloneqq
        \lyapunovrank{\Lnplus[x]}$ satisfying \cref{ineq:f is
          superadditive} and \cref{eq:f of n - 1}. One can derive a
        more-general version of \cref{impl:growth inequality for f},
        namely
        %
        \begin{equation}\label[implication]{impl:growth inequality for f pt2}
          \begin{gathered}
            n - 1 \ge 1 + \dim\of{K} + \delta\\
            \implies\\
            f\of{n - 1} + f\of{1 + \dim\of{K}}
            \ge
            f\of{n - 1 - \delta} + f\of{1 + \dim\of{K} + \delta}.
          \end{gathered}
        \end{equation}

        As before, we consider two cases.

        \begin{pcases}
          \begin{case}[$n_{1} \le \dim\of{K}$]
            This time around we define $\smallestidx$ to be the
            smallest index such that $S\of{\smallestidx} >
            \dim\of{K}$. Then $S\of{k} \le 2\dim\of{K}$, and we need
            $n > \dim\of{K}$ to conclude that $\smallestidx < k$.
            \cref{lem:more conditions} gives us that, and moreover
            lets us apply \cref{impl:growth inequality for f pt2}.
            Then finally, \cref{itm:condition2} lets us conclude that
            $\sum_{i=1}^{k}f\of{n_{i}} < f\of{n} + \lyapunovrank{K}$.
          \end{case}

          \begin{case}[$n_{1} > \dim\of{K}$]
            With $\sum_{i=2}^{k} n_{i} = \dim\of{K} + 1 + \delta$, the
            proof proceeds as before. By assumption we can apply
            \cref{impl:growth inequality for f pt2}, and then
            \cref{itm:condition2} gives $\sum_{i=1}^{k}f\of{n_{i}} <
            f\of{n} + \lyapunovrank{K}$. \qedhere
          \end{case}
          \end{pcases}
      \end{proof}
    \end{lemma}


    \begin{theorem}\label{thm:subproblem reduction}
      Suppose that $K,J$ are symmetric cones and that $n \in \Nn[1]$
      satisfies \cref{itm:condition1,itm:condition2,itm:condition3}.
      Then $J \sim \directsum{\Lnplus}{K}$ if and only if there exists
      a symmetric $J'$ such that $J \cong \directsum{\Lnplus}{J'}$ and
      $J' \sim K$.

      \begin{proof}
        The ``if'' case is obvious and doesn't require any conditions
        on $n$. Likewise when $K = \set{0}$. In the other direction,
        and with $\dim\of{K} \ge 1$, we will show that if $J \sim
        \directsum{\Lnplus}{K}$ does not have an $\Lnplus$ factor,
        then there exists some symmetric cone $H$ having $\dim\of{H} =
        \dim\of{J}$ but $\lyapunovrank{J} \le \lyapunovrank{H} <
        \lyapunovrank{\directsum{\Lnplus}{K}}$. This of course would
        be a contradiction, because $J$ and $\directsum{\Lnplus}{K}$
        share a signature.

        Suppose, for example, that $J \sim \directsum{\Lnplus}{K}$ and
        that $J$ has an $\Hnplus[m]\of{\Cn[1]}$ factor with $m \ge
        3$. We will ``replace'' this factor with a sum of Lorentz
        cones. With \cref{lem:not all smaller factors} in mind, what
        we would like to do is find a set of Lorentz cones each of
        dimension strictly less than $n$ whose dimensions add up to
        the dimension of $\Hnplus[m]\of{\Cn[1]}$ and whose combined
        Lyapunov ranks exceed that of $\Hnplus[m]\of{\Cn[1]}$. Since
        $\Hnplus[m]\of{\Cn[1]}$ is a factor of $J$, and since
        $\dim\of{J} = n + \dim\of{K}$, we can say with certainty that
        %
        \begin{equation*}
          \dim\of{\Hnplus[m]\of{\Cn[1]}} = m^{2} \le n + \dim\of{K} < 2n,
        \end{equation*}
        %
        where the last inequality comes from \cref{lem:more
          conditions}. Thus, $m^{2}/2 < n$, and this tells us exactly
        how big our Lorentz cones can be if we intend to apply
        \cref{lem:not all smaller factors}. There are two cases for
        $m$, each of which leads to a different set of Lorentz cones:
        %
        \begin{equation*}
          I
          \coloneqq
          \begin{cases*}
            \Lnplus[m^{2}/2]
            \oplus
            \Lnplus[m^{2}/2] &
              if $m$ is even,\\
            \Lnplus[\qty{m^{2} - 1}/2]
            \oplus
            \Lnplus[\qty{m^{2} - 1}/2]
            \oplus
            \Lnplus[1]
            & if $m$ is odd.
          \end{cases*}
        \end{equation*}
        %
        In either case, $\dim\of{I} = \dim\of{\Hnplus[m]\of{\Cn[1]}}$,
        and the Lorentz factors of $I$ have dimension strictly less
        than $n$. We would also like
        $\lyapunovrank{\Hnplus[m]\of{\Cn[1]}} < \lyapunovrank{I}$, and
        it is easy to check that $\lyapunovrank{I} -
        \lyapunovrank{\Hnplus[m]\of{\Cn[1]}} > 0$ holds if $m \ge
        4$. For $m = 3$, we may take $I \coloneqq \Lnplus[9]$ instead
        of using the formula, since $9 < n$ by
        \cref{itm:condition3}.

        Now suppose we take the remaining factors of $J$ and group
        them into $\bar{J}$, so that $J \cong
        \directsum{\Hnplus[m]\of{\Cn[1]}}{\bar{J}}$. Then for all $m
        \ge 3$,
        %
        \begin{equation*}
          \lyapunovrank{J}
          =
          \lyapunovrank{\Hnplus[m]\of{\Cn[1]}} + \lyapunovrank{\bar{J}}
          <
          \lyapunovrank{I} + \lyapunovrank{\bar{J}}.
        \end{equation*}
        %
        Define $H \coloneqq \directsum{I}{\bar{J}}$. Does $H$ consist
        only of Lorentz factors? If not, repeat the process above
        until it does. The other two irreducible families
        $\Hnplus[m]\of{\Rn[1]}$ and $\Hnplus[m]\of{\quaternions}$ can
        be handled in a similar manner as $\Hnplus[m]\of{\Cn[1]}$, and
        the sole octonion cone $\Hnplus[3]\of{\octonions}$ can be
        replaced by $I \coloneqq \Lnplus[9] \oplus \Lnplus[9] \oplus
        \Lnplus[9]$.

        After a finite number of steps, we may assume that $H
        \coloneqq \directsum{I}{\bar{J}}$ where $I$ consists only of
        Lorentz factors of dimension less than $n$, and where
        $\bar{J}$ contains whatever Lorentz factors were in $J$ when
        we started\footnote{If $J$ has no non-Lorentz factors, we set $H
        \coloneqq \bar{J} \coloneqq J$ as step zero.}. By \cref{lem:no
          bigger factors pt2}, none of the factors in $\bar{J}$ are
        larger than $\Lnplus$. And if the factors of $\bar{J}$ were
        all smaller than $\Lnplus$, then every factor of $H$ would be
        smaller than $\Lnplus$, and \cref{lem:not all smaller factors}
        would show that $\lyapunovrank{H} <
        \directsum{\Lnplus}{K}$. But that of course would be a
        contradiction, since $\lyapunovrank{J} \le
        \lyapunovrank{H}$. As a result, at least one of the Lorentz
        factors remaining in $\bar{J}$ is $\Lnplus$, and it originated
        from $J$.

        Take one of those $\Lnplus$ factors from within $J$, and write
        $J \cong \directsum{\Lnplus}{J'}$. It should be clear that if
        $J \sim \directsum{\Lnplus}{K}$, then the signatures of $J'$
        and $K$ agree. Moreover, $J'$ and $K$ cannot be isomorphic;
        for if they were, then $J$ and $\directsum{\Lnplus}{K}$ would
        be isomorphic, and they are not.
      \end{proof}
    \end{theorem}

    \begin{corollary}
      If $K$ is a symmetric cone having no symmetric similacra and if
      $n$ satisfies
      \cref{itm:condition1,itm:condition2,itm:condition3}, then
      $\directsum{\Lnplus}{K}$ has no symmetric similacra.
    \end{corollary}

    \cref{itm:condition3} can be weakened a bit so long as we keep the
    other two. By adding \cref{itm:condition1,itm:condition2}, what we
    obtain is $\dim\of{K}\qty{\dim\of{K} - 1} \le 4n - 10$. For this
    to be true we must have $n \ge 3$, and the potential new cases are
    $n \in \set{3,4,5,\ldots,9}$, because $n \ge 10$ is handled by the
    theorem. For each $n$, then, we must then check all $K$ such that
    $\dim\of{K}\qty{\dim\of{K} - 1} \le 4n - 10$. There are a finite
    number of these, and computation shows that $\Hnplus[3]\of{\Rn[1]}
    \sim \directsum{\Lnplus[4]}{\Rnplus[2]}$ is the sole
    counterexample at $n = 4$. As a result, we could get away with
    requiring $n \ge 5$ instead of $n \ge 10$.
  \end{section}

  \section*{Acknowledgments}{
    The authors thank David Sossa and the University of O'Higgins for
    organizing and sponsoring the 2024 Workshop on Variation Analysis
    and Euclidean Jordan Algebras (VAEJA). The results in
    \cref{sec:reducible cones} were proved as a result of the
    workshop.
  }

  \appendix
  \begin{section}{The table of Lyapunov ranks}
    In the absence of any restrictions on $n$, \cref{tab:irreducible
      signatures} would produce incorrect values for $\Lnplus[0]$,
    $\Hnplus[0]\of{\Cn[1]}$, and $\Hnplus[1]\of{\quaternions}$. The
    theorem of Gowda and Tao (on which our table is based) is itself
    based upon a table at the bottom of page 97 in Faraut and
    Korányi~\cite{faraut-koranyi}. The latter states no restrictions
    on $n$, but that can lead to misinterpretation in two ways. First,
    the rows in the table are not necessarily distinct. When $n = 1$
    for example, we have $\Lnplus = \Hnplus\of{\quaternions}$, and
    when that happens, the choice of the $\Hnplus\of{\quaternions}$
    row (as opposed to the $\Lnplus$ row) is what produces the wrong
    Lyapunov rank. Second, the table is simply not valid in trivial
    cases.

    Based on the notation, we find it likely that the Lie algebras in
    Faraut and Korányi's table are derived from Satz~IX.3.3 in Braun
    in Koecher~\cite{braun-koecher}. Braun and Koecher have one case
    for $\Lnplus$ with $n \ge 2$, and then require $n \ge 3$ for
    $\Hnplus\of{\Rn[1]}$, $\Hnplus\of{\Cn[1]}$, and
    $\Hnplus\of{\quaternions}$. As in \cref{thm:irreducible families},
    this ensures that there is no overlap between the families. From
    this we may draw two conclusions about the Faraut and
    Korányi table:

    \begin{enumerate}
      \begin{item}
        When there is ambiguity, the Lorentz cone row should be used.
      \end{item}
      \begin{item}
        A priori, the table is invalid for $n < 2$.
      \end{item}
    \end{enumerate}

    It is interesting to note however that the entries for
    $\Hnplus[2]\of{\Rn[1]} \cong \Lnplus[3]$, $\Hnplus[2]\of{\Cn[1]}
    \cong \Lnplus[4]$, and $\Hnplus[2]\of{\quaternions} \cong
    \Lnplus[6]$ are consistent. In each case we have a choice of rows
    (and should prefer the row for the Lorentz cone), but up to
    isomorphism, either row gives the same Lie algebras. The entire
    table is therefore free of ambiguity for $n \ge 2$, which is how
    we arrive at the $n \ge 2$ condition on \cref{tab:irreducible
      signatures}.

    For the remaining cases $n \in \set{0,1}$, the computations are
    trivial and we find it much simpler to use a separate table
    (\cref{tab:rnplus signature}) than it would be to make
    \cref{tab:irreducible signatures} consistent for all $n \ge 0$.
  \end{section}

  \bibliographystyle{mjo}
  \bibliography{references}
\end{document}
